% concatenated from torelli_T.tex
\documentclass[11pt,a4paper]{article}
\usepackage[T1]{fontenc}
\usepackage[utf8]{inputenc}
\usepackage{lmodern,amsmath,amssymb,amsthm,mathtools,mathrsfs,txfonts}
\usepackage[margin=35mm]{geometry}
\usepackage{microtype,booktabs,array,enumitem,xcolor}
\usepackage[colorlinks=true,linkcolor=black,citecolor=blue!50!black,urlcolor=blue!50!black]{hyperref}
\usepackage[nameinlink,noabbrev]{cleveref}



\setlength{\emergencystretch}{2em}
\setlist{itemsep=3pt,topsep=5pt}
\newcolumntype{P}[1]{>{\raggedright\arraybackslash}p{#1}}
\newtheorem{theorem}{Theorem}[section]
\newtheorem{maintheorem}{Theorem}
\renewcommand{\themaintheorem}{\Alph{maintheorem}}
\crefname{maintheorem}{Theorem}{Theorems}
\Crefname{maintheorem}{Theorem}{Theorems}
\newtheorem{proposition}[theorem]{Proposition}
\newtheorem{lemma}[theorem]{Lemma}
\newtheorem{corollary}[theorem]{Corollary}
\theoremstyle{definition}
\newtheorem{definition}[theorem]{Definition}
\newtheorem{example}[theorem]{Example}
\newtheorem{remark}[theorem]{Remark}

\crefname{theorem}{Theorem}{Theorems}
\crefname{proposition}{Proposition}{Propositions}
\crefname{lemma}{Lemma}{Lemmas}
\crefname{corollary}{Corollary}{Corollaries}
\crefname{definition}{Definition}{Definitions}
\crefname{example}{Example}{Examples}
\crefname{remark}{Remark}{Remarks}
\crefname{section}{Section}{Sections}

\newcommand{\R}{\mathbb R}
\newcommand{\Q}{\mathbb Q}
\newcommand{\Z}{\mathbb Z}
\newcommand{\C}{\mathbb C}
\newcommand{\T}{\mathbb T}
\newcommand{\PP}{\mathbb P}
\newcommand{\HH}{\mathcal H}
\newcommand{\Mod}{\operatorname{Mod}}
\newcommand{\Sp}{\operatorname{Sp}}
\newcommand{\SL}{\operatorname{SL}}
\newcommand{\GL}{\operatorname{GL}}
\newcommand{\Aut}{\operatorname{Aut}}
\newcommand{\Hom}{\operatorname{Hom}}
\newcommand{\tor}{\operatorname{tor}}
\newcommand{\HS}{\operatorname{HS}}
\newcommand{\Ad}{\operatorname{Ad}}
\newcommand{\tr}{\operatorname{tr}}
\newcommand{\im}{\operatorname{im}}
\newcommand{\ab}{\mathrm{ab}}
\newcommand{\rt}{(T)}
\newcommand{\eps}{\varepsilon}
\newcommand{\one}{\mathbf 1}
\newcommand{\Ga}{\mathbf G_a}
\newcommand{\Lie}{\operatorname{Lie}}
\newcommand{\Ru}{R_{\mathrm u}}
\newcommand{\un}{\mathrm{un}}
\title{Property \texorpdfstring{$\rt$}{(T)} and nonlinearity of mapping class group quotients}
\author{Piotr W. Nowak}


\begin{document}
\maketitle
\begin{abstract}
We give a general method for proving Kazhdan's property~\(\rt\)
for quotients \(G/K_{[c+1]}\), where \(G\) is countable, \(K\) is a normal subgroup
and \(K_{[j]}\) denotes its lower central series. The method
combines an affine realization of \(G/[K,K]\), a contraction
argument, and permanence for nilpotent normal subgroups.

We apply it to the Torelli lower-central quotients
\(\Mod(\Sigma_g)/(\mathcal{T}_g)_{[c+1]}\) for \(g\ge3\) and show that they have property~\(\rt\) for every \(c\ge1\).
We also prove
property~\(\rt\) for
\(\Aut(F_3)/(\mathrm{IA}_3)_{[c+1]}\) for every \(c\ge1\), relating
these groups to the tame nilpotent images studied by
Lubotzky and Pak. 

For the Torelli lower-central quotients with \(g\ge3\) and
\(c\ge2\), every finite-dimensional complex representation
has infinite kernel, and these quotients are not linear over
any field.
\end{abstract}


\section{Introduction}

Kazhdan's property~\(\rt\) is a rigidity property of unitary
representations: almost invariant vectors force the existence of
a nonzero invariant vector. For countable discrete groups it is
equivalent to the fixed point property for affine isometric actions
on Hilbert spaces. The examples coming from higher-rank Lie groups
and their lattices have been an important source of both results and
questions in geometric group theory. See \cite{BHV} for an overview of property \(\rt\) and its applications.

Two natural families of interest in this context are automorphism groups of free
groups and mapping class groups. 
For the free group \(F_n\) of rank \(n\), the group \(\Aut(F_n)\)
has property~\(\rt\) for every \(n\ge4\): see \cite{KalubaNowakOzawa} for the case \(n=5\), \cite{KalubaKielakNowak} for the case \(n\ge 5\) and \cite{NitscheAutFour} for \(n=4\).
In rank three, \(\Aut(F_3)\) does not have property~\(\rt\),
as follows from work of McCool \cite{McCool} and Grunewald and Lubotzky
\cite{GrunewaldLubotzky}.
For a closed connected oriented surface \(\Sigma_g\) of genus
\(g\ge3\), it remains open whether \(\Mod(\Sigma_g)\)
has property~\(\rt\).

In this paper we give a general method for proving property~\(\rt\)
for quotients obtained by replacing a normal subgroup with its
nilpotent quotients. The method applies to both families above.
It uses the action on the abelianization of the normal subgroup,
together with an affine realization of the resulting extension.
Once this first quotient has property~\(\rt\), a permanence theorem
for nilpotent normal subgroups gives all subsequent stages.
The first of our main results isolates the hypotheses needed for our
principal applications.

For a group \(K\), denote its lower central series by
\[
 K_{[1]}=K,\qquad K_{[j+1]}=[K,K_{[j]}].
\]

\begin{maintheorem}\label{thm:strategy}
Let \(G\) be a finitely generated group and \(K\subseteq G\)
a normal subgroup with finitely generated abelianization. Put
\[
 A=\bigl(K/[K,K]\bigr)/\tor\bigl(K/[K,K]\bigr),\qquad
 V=A\otimes_\Z\R,\qquad \Gamma=G/K.
\]
Suppose that:
\begin{enumerate}[label=\textup{(\roman*)}]
 \item \(\Gamma\) is a lattice in a second-countable unimodular
 locally compact group \(S\) with property~\(\rt\);
 \item the conjugation action of \(\Gamma\) on \(V\) extends to
 a continuous representation \(\rho:S\to\GL(V)\) satisfying
 \(\lvert\det\rho(s)\rvert=1\) for every \(s\in S\);
 \item \(V\) is spanned by vectors \(v\) for which there is a
 sequence \((s_j)\) in \(S\) such that \(\rho(s_j)v\to0\);
 \item some element \(z\) of the center of \(\Gamma\) fixes no
 nonzero vector of \(V\).
\end{enumerate}
Then \(G/K_{[c+1]}\) has property~\(\rt\) for every \(c\ge1\).
More generally, every quotient of \(G\) in which the image of \(K\)
is nilpotent has property~\(\rt\).
\end{maintheorem}

In both applications the central element in~(iv) acts as \(-I\).
A lift of this element provides affine coordinates for
\(G/[K,K]\), after removing the finite torsion in its abelian
kernel. No splitting of the extension is required.
The contraction assumption~(iii) gives property~\(\rt\) for
\(V\rtimes S\), and the affine realization identifies the
discrete group as a lattice in it. Finally, the theorem of
Chatterji, Witte Morris and Shah \cite{CWMS} implies that
\[
 G/K_{[c+1]}\text{ has property }\rt
 \quad\Longleftrightarrow\quad
 G/[K,K]\text{ has property }\rt,\qquad c\ge1.
\]
Thus no separate analysis of the higher lower-central quotients
of \(K\) is needed. We prove a more general affine criterion in
\cref{thm:generalcriterion}, from which \cref{thm:strategy} follows.

We first apply this method to mapping class groups. Denote by
\(H_g=H_1(\Sigma_g;\Z)\) the first homology of the surface and by
\[
 \mathcal{T}_g
 =\ker\bigl(\Mod(\Sigma_g)\longrightarrow\Sp_{2g}(\Z)\bigr)
\]
the Torelli group. For \(c\ge1\), set
\[
 \mathcal{Q}_{g,c}=\Mod(\Sigma_g)/(\mathcal{T}_g)_{[c+1]},
 \qquad
 \mathcal{N}_{g,c}=\mathcal{T}_g/(\mathcal{T}_g)_{[c+1]}.
\]
The quotient \(\mathcal{Q}_{g,c}/\mathcal{N}_{g,c}\) is
\(\Sp_{2g}(\Z)\).
Johnson's abelianization theorem \cite{JohnsonIII} identifies
the torsion-free abelianization of \(\mathcal{T}_g\) with
\[
 A_g=\Lambda^3H_g/(H_g\wedge\Omega),
\]
where \(\Omega\) is the inverse symplectic form.
Its action extends to \(\Sp_{2g}(\R)\); the central element
\(-I\) acts as \(-1\), and a diagonal subgroup contracts a
spanning set. These facts verify the hypotheses of
\cref{thm:strategy} and give our first application.

\begin{maintheorem}\label{thm:main}
For every \(g\ge3\) and every integer \(c\ge1\), the group
\(\mathcal{Q}_{g,c}\) is finitely generated, infinite and has
Kazhdan's property~\(\rt\). Its normal subgroup \(\mathcal{N}_{g,c}\)
is finitely generated and nilpotent, has infinite abelianization,
and satisfies \(\mathcal{Q}_{g,c}/\mathcal{N}_{g,c}\cong\Sp_{2g}(\Z)\).
\end{maintheorem}

The genus assumption is essential. For every \(c\ge1\), the
group \(\mathcal{Q}_{2,c}\) has a finite-index subgroup surjecting
onto \(\Z\), and hence does not have property~\(\rt\).
The first-stage affine realization for \(g\ge3\) already appears
in the work of Hain and Matsumoto \cite[Lemma~6.3]{HainMatsumoto}.
In \cref{sec:groups} we explain how it enters the argument.

We then consider the case of \(\Aut(F_3)\), the automorphism group of the free group on \(3\) generators.
For \(\Aut(F_3)\), the relevant normal subgroup is
\[
 \operatorname{IA}_3
 =\ker\bigl(\Aut(F_3)\longrightarrow\GL_3(\Z)\bigr).
\]
Denote \(H_3=H_1(F_3;\Z)\) and \(H_3^*=\Hom(H_3,\Z)\).
The first Johnson homomorphism (see \cite{Johnson1980} for its Torelli-group
version) identifies the abelianization of
\(\operatorname{IA}_3\) with
\(H_3^*\otimes\Lambda^2H_3\), see \cite{SatohSurvey} for details.
Here the arithmetic quotient is \(\GL_3(\Z)\), the ambient group
is
\(\{s\in\GL_3(\R):\lvert\det s\rvert=1\}\), and again
\(-I\) acts as \(-1\) on the abelianized kernel.
The contraction condition is verified by a diagonal subgroup.
With
\[
 \mathcal{P}_{3,c}
 =\Aut(F_3)/(\operatorname{IA}_3)_{[c+1]},
\]
we obtain the following application.

\begin{maintheorem}\label{thm:autquotients}
For every integer \(c\ge1\), the group
\(\mathcal{P}_{3,c}\) is finitely generated, infinite and has
property~\(\rt\). The normal subgroup
\(\operatorname{IA}_3/(\operatorname{IA}_3)_{[c+1]}\) is
finitely generated and nilpotent, and its abelianization is the
infinite free abelian group
\[
 H_3^*\otimes\bigwedge\nolimits^2 H_3.
\]
In particular, every quotient of \(\Aut(F_3)\) in which
\(\operatorname{IA}_3\) has nilpotent image has property~\(\rt\).
\end{maintheorem}

\Cref{thm:autquotients} is stated for rank \(n=3\), although the same argument also applies for \(n\ge4\).
For \(n\ge4\), the
analogous property~\(\rt\) assertion already follows from
property~\(\rt\) of \(\Aut(F_n)\) and permanence under quotients.
Recall that \(\operatorname{Aut}(F_3)\) does not have
property~\(\rt\) \cite{McCool,GrunewaldLubotzky}, although every \(\mathcal{P}_{3,c}\) does.
In \cref{sec:freeautomorphisms} we discuss the relation with
Lubotzky and Pak's earlier results on tame automorphisms of free
nilpotent groups \cite{LubotzkyPak}, which also give the first
stage of this application.

An important feature of the argument in Theorem \ref{thm:strategy} is that the same method
applies to both families, \(\Mod(\Sigma_g)\) and \(\Aut(F_n)\). The homology modules and arithmetic
groups differ, but the passage from the action on first homology
to an affine lattice, and then to every nilpotent stage, is the
same. 

We finally consider linear representations of the mapping-class
quotients \(\mathcal{Q}_{g,c}\). The usual symplectic quotient is
linear, and the first-stage quotient, after removal of its finite
torsion kernel, has the affine realization used above. Starting
at \(c=2\), a different feature of the Torelli group survives.
Hain's comparison of unipotent and relative completions
\cite{Hain} gives a rational central line which remains nonzero
in the unipotent completions of these nilpotent quotients but is killed in the relative
completion of the mapping class group. Together with arithmetic
superrigidity, this gives a stronger obstruction than the
absence of a faithful representation.

\begin{maintheorem}\label{thm:nonlinearity}
Let \(g\ge3\) and \(c\ge2\). Every homomorphism
\[
 \mathcal{Q}_{g,c}\longrightarrow\GL_d(\C)
\]
has infinite kernel. In particular, \(\mathcal{Q}_{g,c}\) is not linear
over \(\C\).
\end{maintheorem}

The proof uses superrigidity to place a hypothetical finite-kernel representation into a form to which the universal property of relative completion applies.
The central line must then be killed, and rational Malcev theory
produces an infinite subgroup of the discrete kernel.
A separate argument excludes faithful representations in
positive characteristic. Consequently \(\mathcal{Q}_{g,c}\)
is not linear over any field when \(g\ge3\) and \(c\ge2\), 
see \cref{cor:charzero,cor:allfields}.
These groups thus differ from the familiar arithmetic quotients
not only in their kernels but also in their linear representation
theory. It is not clear  if the groups 
\(\mathcal{P}_{3,c}\) are linear.

Hull's common-quotient theorem \cite[Corollary~1.6]{Hull} and
SQ-universality \cite[Theorem~8.1(a)]{DGO} also give, for
\(g\ge3\), a quotient of \(\Mod(\Sigma_g)\) with property~\(\rt\)
containing \((\Q,+)\). This quotient is finitely generated but
not residually finite, and hence is nonlinear over every field
by Malcev's theorem \cite{Nica}.
The advantage of the construction in Theorem \ref{thm:main} is its
canonically defined tower, with finitely generated nilpotent
kernels over \(\Sp_{2g}(\Z)\). It preserves the standard homological
quotient and identifies the obstruction to linearity within
the lower central structure of the Torelli group.

The paper is organized as follows. In \cref{sec:general} we
develop the general strategy. We apply it to mapping class
groups in \cref{sec:groups} and to automorphism groups of free
groups in \cref{sec:freeautomorphisms}. The proof of
\cref{thm:nonlinearity} is given in \cref{sec:nonlinearity}.
Further applications are discussed in \cref{sec:finalremarks}.

\tableofcontents

\section{A general strategy for nilpotent-kernel quotients}
\label{sec:general}
\label{sec:permanence}

In this section we develop a general argument for proving property~\(\rt\).
It applies to an arbitrary
countable group with a normal subgroup and consists of two parts: algebraic and analytic. 
The algebraic part reduces property \(\rt\) for
an entire family of nilpotent-kernel quotients to proving it for a single quotient,
obtained by abelianizing the kernel. The analytic part  then gives a
sufficient condition for this first quotient to have
property~\(\rt\): an affine realization over a lattice in a
locally compact group with property~\(\rt\), together with
contraction of the translation directions.
The hypotheses concerning lattices,
modules and contraction will be checked separately in the applications.

\subsection{Conventions and elementary permanence properties}

We use the commutator convention \([x,y]=xyx^{-1}y^{-1}\).
For subgroups, brackets denote the subgroup generated by the
corresponding commutators. The lower central series of a group
\(G\) is
\[
 G_{[1]}=G,\qquad G_{[j+1]}=[G,G_{[j]}]\quad(j\ge1).
\]
The group is \emph{nilpotent of class at most \(c\)} if
\(G_{[c+1]}=1\). The abelianization is \(G/[G,G]\).
If \(H\subseteq G\) is a normal subgroup, all its lower-central subgroups are
characteristic in \(H\), hence normal in \(G\).
For a finitely generated abelian group \(B\), write
\(\tor(B)\) for its finite torsion subgroup.

All locally compact groups below are Hausdorff. A unitary
representation \(\pi\) of such a group \(G\) on a complex
Hilbert space \(\HH\) is assumed to be strongly continuous: for each
\(\xi\in\HH\), the map \(g\mapsto\pi(g)\xi\) is norm
continuous. It has \emph{almost invariant unit vectors} if, for
every compact subset \(C\subset G\) and every \(\eps>0\),
there is \(\xi\in\HH\) with \(\|\xi\|=1\) and
\[
 \sup_{s\in C}\|\pi(s)\xi-\xi\|<\eps.
\]
For a subgroup \(H\subseteq G\), denote
\[
 \HH^H=\{\xi\in\HH:\pi(h)\xi=\xi\text{ for every }h\in H\}.
\]
The group has \emph{Kazhdan's property~\(\rt\)} if every unitary representation \(\pi\) of \(G\) with almost invariant vectors satisfies \(\HH^G\neq 0\).

The pair \((G,H)\) has \emph{relative property~\(\rt\)} if every unitary representation \(\pi\) of \(G\) with almost invariant vectors has \(\HH^H\ne0\).

We shall use the following standard permanence properties for
countable discrete groups. Property~\(\rt\) passes to quotient
groups \cite[Theorem~1.3.4]{BHV}. If \(H\subseteq G\) is a
subgroup of finite index, then \(G\) has property~\(\rt\) if
and only if \(H\) does \cite[Theorem~1.7.1]{BHV}. Property~\(\rt\)
is also invariant under passage through finite normal kernels:
if \(F\subseteq G\) is a finite normal subgroup, then \(G\)
has property~\(\rt\) if and only if \(G/F\) does
\cite[Propositions~1.1.5 and~1.7.6]{BHV}.

\begin{lemma}
\label{lem:normalextension}
Let \(N\subseteq G\) be a normal subgroup of a countable
discrete group. The following conditions are equivalent:
\begin{enumerate}[label=\textup{(\roman*)},beginpenalty=10000,midpenalty=10000]
 \item \(G\) has property~\(\rt\);
 \item \((G,N)\) has relative property~\(\rt\) and \(G/N\)
 has property~\(\rt\).
\end{enumerate}
\end{lemma}
\begin{proof}
Suppose~(i) holds. The relative assertion in~(ii) follows
from the definition, and the assertion about \(G/N\) follows
from quotient permanence.

Conversely, assume~(ii), and let \(\pi\) have almost invariant unit vectors.
Choose finite subsets \(C_n\) increasing to \(G\), and unit
vectors \(\xi_n\) such that
\[
 \sup_{s\in C_n}\|\pi(s)\xi_n-\xi_n\|<1/n.
\]
In particular, for every fixed \(s\in G\),
\(\|\pi(s)\xi_n-\xi_n\|\to0\).
The orthogonal projection
\(P:\HH\to\HH^N\) commutes with every \(\pi(s)\).
We claim that \(\|(1-P)\xi_n\|\to0\).
Otherwise, some subsequence satisfies
\(\|(1-P)\xi_{n_j}\|\ge\delta>0\).
The  normalized vectors \( \eta_j=(1-P)\xi_{n_j}/\|(1-P)\xi_{n_j}\| \in(\HH^N)^\perp \)
are unit vectors, and we have
\[
 \begin{aligned}
 \|\pi(s)\eta_j-\eta_j\| \le\delta^{-1}\|\pi(s)\xi_{n_j}-\xi_{n_j}\|.
 \end{aligned}
\]
For every finite set of \(s\), the right-hand side tends
uniformly to zero. Hence the representation on
\((\HH^N)^\perp\) has almost invariant unit vectors.
Since it has no nonzero \(N\)-fixed vector, this contradicts
relative property~\(\rt\).

We have  \(\|P\xi_n\|^2=1-\|(1-P)\xi_n\|^2\to1\).
For all sufficiently large \(n\), the vectors
\(P\xi_n/\|P\xi_n\|\) are defined, and the same estimate,
with \(P\) in place of \(1-P\), shows that they are almost
invariant. The representation on \(\HH^N\) factors through
\(G/N\). Property~\(\rt\) of this quotient supplies a
guarantees the existence of a non-zero vector, which is fixed by all of \(G\).
\end{proof}

\subsection{Nilpotent normal subgroups and quotient towers}

The following result is the starting point for our investigation.
\begin{theorem}[Chatterji--Witte Morris--Shah]
\label{thm:cwms}
Let \(N\subseteq G\) be a nilpotent normal subgroup of a countable
discrete group. The following conditions are equivalent:
\begin{enumerate}[label=\textup{(\roman*)},beginpenalty=10000,midpenalty=10000]
 \item \((G,N)\) has relative property~\(\rt\);
 \item \(\bigl(G/[N,N],\,N/[N,N]\bigr)\) has relative
 property~\(\rt\).
\end{enumerate}
\end{theorem}

This result is \cite[Theorem~1.2]{CWMS} in the special case of discrete groups. 

\begin{theorem}
\label{thm:nilpotentpermanence}
Let \(N\subseteq G\) be a nilpotent normal subgroup, with \(G\)
countable and discrete. The following conditions are equivalent:
\begin{enumerate}[label=\textup{(\roman*)},beginpenalty=10000,midpenalty=10000]
 \item \(G\) has property~\(\rt\);
 \item \(G/[N,N]\) has property~\(\rt\).
\end{enumerate}
\end{theorem}
\begin{proof}
The implication \(\textup{(i)}\Rightarrow\textup{(ii)}\) follows since property~\(\rt\) is preserved by quotients. 
Conversely, assume~(ii). Then \(G/[N,N]\) has relative property~\(\rt\)
with respect to every subgroup, in particular \(N/[N,N]\).
By \cref{thm:cwms}, \((G,N)\) has relative property~\(\rt\).
Also \(G/N\), being a quotient of \(G/[N,N]\), has
property~\(\rt\). The claim follows from \cref{lem:normalextension}.
\end{proof}

\begin{remark}
\label{rem:nilpotenceessential}
The hypothesis cannot be replaced by virtual nilpotence or
solvability. Indeed, let
\[
 G=N=D_\infty
 =\langle r,s\mid s^2=1,\ srs^{-1}=r^{-1}\rangle
 \cong\Z\rtimes C_2.
\]
Then \([N,N]=\langle r^2\rangle\) and
\(G/[N,N]\cong C_2\times C_2\) is finite, hence has
property~\(\rt\). However, \(\langle r\rangle\cong\Z\)
has index two in \(G\), so \(G\) does not have
property~\(\rt\).
This subgroup \(N\) is virtually cyclic, virtually nilpotent,
metabelian and polycyclic. Thus the same example excludes each
of these properties, as well as solvability. 
\end{remark}

\begin{remark}
\label{rem:nilpotentextensions}
The conclusion of \cref{thm:nilpotentpermanence} remains valid
if there is a normal subgroup \(K\subseteq G\), with \(K\subseteq N\), such that
\((G,K)\) has relative property~\(\rt\) and \(N/K\) is
nilpotent. To see the reverse implication, observe that
\[
 (G/K)/[N/K,N/K]\cong G/(K[N,N])
\]
is a quotient of \(G/[N,N]\). Applying \cref{thm:nilpotentpermanence}
to the nilpotent normal subgroup \(N/K\subseteq G/K\) gives
property~\(\rt\) for \(G/K\). Relative property~\(\rt\) for \(K\)
then gives property~\(\rt\) for \(G\).
It suffices, in particular, that \(K\) itself has
property~\(\rt\) and \(N/K\) is nilpotent.

Another special case is when \(N\) is \emph{finite-by-nilpotent}:
there is a finite normal subgroup \(F\subseteq N\) with
\(N/F\) nilpotent. If its nilpotency class is at most \(c\),
then \(K=N_{[c+1]}\subseteq F\) is finite and characteristic
in \(N\), hence normal in \(G\). Apply the preceding argument
to this \(K\). The initially chosen \(F\) need not be normal
in \(G\). The order of the words in ``finite-by-nilpotent''
matters: the opposite order is ruled out by the dihedral example.
\end{remark}

The next statement shows that property~\(\rt\) for the first quotient determines it for the whole tower of quotients.
\begin{theorem}
\label{thm:quotienttower}
Let \(K\subseteq G\) be a normal subgroup, where \(G\) is countable
and discrete. For every \(c\ge1\), the following conditions are equivalent:
\begin{enumerate}[label=\textup{(\roman*)},beginpenalty=10000,midpenalty=10000]
 \item \(G/K_{[c+1]}\) has property~\(\rt\);
 \item \(G/[K,K]\) has property~\(\rt\).
\end{enumerate}
If these conditions hold, every quotient
of \(G\) in which the image of \(K\) is nilpotent has
property~\(\rt\). The same conclusion holds if that image
is finite-by-nilpotent.
\end{theorem}
\begin{proof}
The subgroup \(K/K_{[c+1]}\) is nilpotent and normal in
\(G/K_{[c+1]}\), and
\[
 \begin{aligned}
 [K/K_{[c+1]},K/K_{[c+1]}]&=[K,K]/K_{[c+1]},\\
 (G/K_{[c+1]})/\bigl([K,K]/K_{[c+1]}\bigr)&\cong G/[K,K].
 \end{aligned}
\]
The equivalence is therefore \cref{thm:nilpotentpermanence}.

More generally, let \(p:G\rightarrow E\) be a surjective homomorphism, and write
\(N=p(K)\). There is a surjection
\[
 G/[K,K]\rightarrow E/[N,N],
\]
because \(p([K,K])=[N,N]\). Thus \(E/[N,N]\) has
property~\(\rt\). If \(N\) is nilpotent, apply
\cref{thm:nilpotentpermanence}; if it is finite-by-nilpotent,
apply \cref{rem:nilpotentextensions}.
\end{proof}

Finite generation is not needed for this theorem. When \(G\)
is finitely generated its quotients are clearly finitely
generated as well. No  assumption about finite generation of \(K\)
is implicit here.

\subsection{Contracting affine groups}
\label{sec:generalaffine}

We will now describe a contraction argument, that will be crucial for proving property \(\rt\).
Let \(V\) be a finite-dimensional real vector space and
\(\rho:S\to\GL(V)\) a continuous representation of a
locally compact group. The \emph{semidirect product}
\(V\rtimes_\rho S\) has underlying space \(V\times S\) and group operation defined by
\[
 (v,s)(w,t)=(v+\rho(s)w,st),\qquad
 (v,s)^{-1}=(-\rho(s^{-1})v,s^{-1}).
\]
We omit the subscript when the action is understood. The vector
subgroup \(V\times\{1\}\) consists of translations, and
\(\{0\}\times S\) is identified with \(S\).

\begin{lemma}
\label{lem:contraction}
Given a continuous unitary representation \(\pi\) of
\(V\rtimes S\), suppose that \(\xi\) is fixed by \(S\).
If \(v\in V\) and some sequence \(s_n\in S\) satisfies
\(\rho(s_n)v\to0\), then \(\xi\) is fixed by the
translation \((v,1)\).
\end{lemma}
\begin{proof}
Conjugation by \((0,s_n)\) sends \((v,1)\) to
\((\rho(s_n)v,1)\). Since \(\xi\) is fixed by
\((0,s_n)\) and its inverse, unitarity gives
\[
 \|\pi(v,1)\xi-\xi\|
 =\|\pi(\rho(s_n)v,1)\xi-\xi\|\longrightarrow0.
\]
The convergence is strong continuity at the identity of the
translation subgroup. The expression on the left is independent
of \(n\), and is therefore zero.
\end{proof}

The above lemma is an elementary form of the Mautner phenomenon,
see \cite[Lemma~1.4.8]{BHV}.

\begin{proposition}
\label{prop:affineT}
Suppose \(S\) has property~\(\rt\). If \(V\) is spanned
by vectors each of which can be contracted to zero, in the sense of Lemma \ref{lem:contraction}, by a sequence
in \(S\), then \(V\rtimes S\) has property~\(\rt\).
\end{proposition}
\begin{proof}
Let a unitary representation \(\pi\) have almost invariant unit vectors. Its restriction
to \(S\) has almost invariant vectors, because a compact subset
of \(S\) is compact in \(V\rtimes S\). Property~\(\rt\)
of \(S\) gives a nonzero vector \(\xi\) fixed by \(S\).
By \cref{lem:contraction}, it is fixed by every contractible
translation. A sequence contracting \(v\) also contracts every
real multiple of \(v\). Every vector in \(V\) is a finite
sum of these multiples, and translations add under multiplication.
Thus \(\xi\) is fixed by all of \(V\), and hence by
\(V\rtimes S\).
\end{proof}


We next explain the passage to discrete groups. 
A \emph{left Haar measure} on a
locally compact group is a nonzero left-translation-invariant
Radon measure; it exists and is unique up to a positive scalar.
A group is \emph{unimodular} if its left Haar measure is also
right invariant. A discrete subgroup \(\Gamma\) of a locally
compact group \(S\) is a \emph{lattice} if \(S/\Gamma\)
admits a nonzero finite \(S\)-invariant Radon measure.
Here \(S/\Gamma\) denotes cosets \(s\Gamma\), obtained
by the right multiplication action of \(\Gamma\) on \(S\).
For a second-countable unimodular group, the lattice condition
can equivalently be expressed by a Borel fundamental domain
of positive finite Haar measure, see
\cite[Sections~A.3 and~B.1--B.2]{BHV}.

A \emph{full lattice} in \(V\) is a subgroup
\(A=\Z b_1+\cdots+\Z b_d\), where \(b_1,\ldots,b_d\)
is a real basis of \(V\). Its quotient \(V/A\) is a compact
torus. We allow \(d=0\), in which case both \(A\) and
\(V\) are zero and their fundamental domain is a point of
measure one.
The next statement gives sufficient conditions for a discrete subgroup of the affine group to be a lattice.
\begin{proposition}
\label{prop:affinelattice}
Let \(S\) be a second-countable unimodular locally compact
group, \(V\) a finite-dimensional real vector space, and
\(\rho:S\to\GL(V)\) a continuous representation preserving
Lebesgue measure, that is,
\[
 |\det\rho(s)|=1\qquad(s\in S).
\]
Let \(p:V\rtimes S\to S\) be the coordinate projection,
and identify \(V\) with the translation subgroup
\(V\times\{1\}\). Suppose that a subgroup
\(\Delta\subseteq V\rtimes S\) satisfies the following:
\begin{enumerate}[label=\textup{(\roman*)},beginpenalty=10000,midpenalty=10000]
 \item \(\Gamma=p(\Delta)\) is a lattice in \(S\);
 \item \(A=\Delta\cap V\) is a full lattice in \(V\).
\end{enumerate}
Then \(\Delta\) is a lattice in \(V\rtimes S\).
In particular, \(A\rtimes\Gamma\) is a
lattice whenever \(\Gamma\) is a lattice in \(S\) preserving
a full lattice \(A\) in \(V\).

If, moreover, \(S\) has property~\(\rt\) and \(V\) is
spanned by vectors each of which can be contracted to zero
by a sequence in \(S\), then both \(V\rtimes S\) and
\(\Delta\) have property~\(\rt\).
\end{proposition}
\begin{proof}
Write \(L=V\rtimes S\). Discreteness of \(\Gamma\) and \(A\)
implies that a sufficiently small product neighborhood of the
identity meets \(\Delta\) only at the identity. Hence
\(\Delta\) is a closed discrete subgroup.
Conjugation in \(\Delta\) also gives \(\rho(\Gamma)A=A\).

Let \(\nu\) be Haar measure on \(S\). Choose a half-open
fundamental parallelepiped \(P\) for \(A\), and normalize
Lebesgue measure \(\lambda\) by \(\lambda(P)=1\).
The determinant assumption and unimodularity of \(S\)
make \(\mu=\lambda\times\nu\) a left and right Haar measure
on \(L\).

Choose a Borel fundamental domain \(\mathcal D\) for
\(S/\Gamma\), with \(0<\nu(\mathcal D)<\infty\), and put
\[
 \mathcal F=\{(\rho(s)u,s):s\in\mathcal D,\ u\in P\}.
\]
This is a Borel fundamental domain for \(L/\Delta\).
Indeed, choose a lift \((b_\gamma,\gamma)\in\Delta\) for
each \(\gamma\in\Gamma\). Given \((v,t)\in L\), write
uniquely \(t=s\gamma\), with \(s\in\mathcal D\), and
\(\rho(s)^{-1}v-b_\gamma=u+a\), with \(u\in P\), \(a\in A\).
Then
\[
 (v,t)=(\rho(s)u,s)(a+b_\gamma,\gamma),
 \qquad (a+b_\gamma,\gamma)\in\Delta.
\]
These unique decompositions give exactly one representative
in \(\mathcal F\) for each coset, without requiring a
homomorphic choice of lifts.

Volume preservation now gives
\[
 0<\mu(\mathcal F)
 =\int_{\mathcal D}\lambda(\rho(s)P)\,d\nu(s)
 =\nu(\mathcal D)<\infty.
\]
Thus \(\Delta\) is a lattice by
\cite[Proposition~B.2.4]{BHV}. Equivalently, restricting
\(\mu\) to \(\mathcal F\) and passing to \(L/\Delta\)
gives its finite invariant quotient measure, see
\cite[Corollary~B.1.7 and Proposition~B.2.4]{BHV}.

Finally, under the additional hypotheses, \(L\) has
property~\(\rt\) by \cref{prop:affineT}, and so does its
lattice \(\Delta\) by \cite[Theorem~1.7.1]{BHV}.
\end{proof}

\subsection{The main criterion for property \(\rt\)}

Consider an exact sequence
\begin{equation}\label{eq:extensionaction}
 1\longrightarrow A\longrightarrow E
 \xrightarrow{p}\Gamma\longrightarrow1
\end{equation}
with \(A\) free abelian of finite rank. We will use additive notation for the group operation in \(A\). 
Conjugation on \(A\)
factors through \(\Gamma\), since conjugation by \(A\)
is trivial on \(A\), inducing a map
\[ \rho: \Gamma \to \operatorname{Aut}(A). \]
A map \(b:E\to V\),
where \(V=A\otimes\R\), is a \emph{crossed homomorphism}
if
\[
 b(xy)=b(x)+\rho(p(x))b(y).
\]
This identity says exactly that
\(x\mapsto(b(x),p(x))\) is a homomorphism into
\(V\rtimes\Gamma\). If \(b|_A\) is the inclusion of
\(A\) in \(V\), that homomorphism is injective: an element
of its kernel lies in \(A\), where its first coordinate is
the element itself.

The following elementary construction often provides such a map.
Its usefulness lies in the fact that an actual splitting of \(E\) is not required.

\begin{proposition}\label{prop:centralaffine}
Consider the extension \eqref{eq:extensionaction} and suppose there is \(z\in Z(\Gamma)\)
such that
\[
 D=I-\rho(z):V\longrightarrow V
\]
is invertible. Then there is an injective homomorphism
\[
 \phi:E\longrightarrow A\rtimes\Gamma
\]
such that
\[
 \operatorname{pr}_{\Gamma}\circ\phi=p,
 \qquad \phi(a)=(Da,1)\quad(a\in A),
\]
where \(\operatorname{pr}_{\Gamma}:A\rtimes\Gamma\to\Gamma\)
is the coordinate projection \((a,\gamma)\mapsto\gamma\).
The image of \(\phi\) has index
\[
 [A:D(A)]=|\det D|.
\]
There is also an injective homomorphism into
\(V\rtimes\Gamma\) which restricts to the identity on \(A\).
In particular, if \(z\) acts as \(-1\) on \(A\), the first
index is \(2^{\operatorname{rank}A}\).
\end{proposition}
\begin{proof}
Choose a lift \(u\in E\) of \(z\), and put
\(d(x)=[x,u]\). Since \(z\) is central in \(\Gamma\),
the image of this commutator in \(\Gamma\) is trivial;
hence \(d(x)\in A\). The identity
\[
 [xy,u]=x[y,u]x^{-1}[x,u]
\]
and commutativity in \(A\) imply
\[
 d(xy)=d(x)+\rho(p(x))d(y).
\]
For \(a\in A\), conjugation by \(u\) is \(\rho(z)\), so
\[
 d(a)=a-\rho(z)a=Da.
\]
It follows that \(\phi(x)=(d(x),p(x))\) is a homomorphism.
If \(\phi(x)=1\), then \(p(x)=1\), so \(x=a\in A\),
and \(Da=0\). Invertibility of \(D\) gives \(a=0\).

The image of \(\phi\) projects onto \(\Gamma\), and its
intersection with the translation subgroup is \(D(A)\).
Every coset of this image in \(A\rtimes\Gamma\) has a
representative in \(A\): multiply a given element by the
inverse of an image element having the same second coordinate.
Two translation representatives give the same coset exactly
when their difference belongs to \(D(A)\).
Therefore the index is \([A:D(A)]\).
The operator \(D\) has an integral matrix in a basis of \(A\)
and nonzero determinant. The usual lattice index formula gives
\([A:D(A)]=|\det D|\).

Because \(z\) is central, \(D\) commutes with every
\(\rho(\gamma)\), as does \(D^{-1}\). Thus
\[
 b(x)=D^{-1}d(x)
\]
is a crossed homomorphism into \(V\), with \(b(a)=a\).
The resulting normalized affine embedding has translation
kernel exactly \(A\). If \(\rho(z)=-I\), then \(D=2I\),
which gives the final assertion.
\end{proof}

The normalized coordinates may lie in \(D^{-1}A\) rather
than \(A\). In the case \(D=2I\), they can be half-integral.

We can now combine the above results into a strategy for proving property~\(\rt\). 
Theorem \ref{thm:strategy} stated in the introduction is a slightly specialized case of the following theorem.
\begin{theorem}
\label{thm:generalcriterion}
Let \(G\) be a countable discrete group and \(K\subseteq G\)
a normal subgroup. Suppose \(B=K/[K,K]\) is finitely generated abelian.
Denote
\[
 A=B/\tor(B),\quad V=A\otimes_\Z\R,
 \quad\Gamma=G/K,
 \quad E=(G/[K,K])/\tor(B).
\]
The subgroup \(\tor(B)\) is normal in \(G/[K,K]\), and there is
an exact sequence \[1\to A\to E\to\Gamma\to1.\]
Assume the following:
\begin{enumerate}[label=\textup{(\roman*)}]
 \item \(\Gamma\) is a lattice in a second-countable
 unimodular locally compact group \(S\) with property~\(\rt\).
 \item The action of \(\Gamma\) on \(A\), extended linearly
 to \(V\), is the restriction of a continuous representation
 \(\rho:S\to\GL(V)\) satisfying
 \(|\det\rho(s)|=1\) for every \(s\in S\).
 \item The vector space \(V\) is spanned by vectors each of
 which can be contracted to zero by a sequence in \(S\).
 \item There is a crossed homomorphism \(b:E\to V\) whose
 restriction to \(A\) is its natural inclusion in \(V\).
\end{enumerate}
Then \(G/[K,K]\) has property~\(\rt\). Consequently
\(G/K_{[c+1]}\) has property~\(\rt\) for every \(c\ge1\),
as does every quotient of \(G\) with nilpotent or
finite-by-nilpotent image of \(K\).

In condition \textup{(iv)}, it suffices that there be
\(z\in Z(\Gamma)\) for which \(I-\rho(z)\) is invertible
on \(V\). In particular, it suffices that a central element
act as \(-1\).
\end{theorem}
\begin{proof}
Conjugation preserves torsion in the abelian normal subgroup
\(B\), so \(\tor(B)\) is finite and normal as asserted.
The crossed homomorphism in~(iv) embeds \(E\) into
\(V\rtimes S\), with projection \(\Gamma\) and translation
kernel exactly \(A\). Conditions~(i)--(iii) and
\cref{prop:affinelattice} therefore give property~\(\rt\)
for \(E\). The natural surjection
\(G/[K,K]\to E\) has finite kernel \(\tor(B)\),
so permanence under finite normal kernels gives property~\(\rt\) for
\(G/[K,K]\). Apply \cref{thm:quotienttower} for the remaining
statements. The sufficient condition involving a central element is
\cref{prop:centralaffine}.
\end{proof}

For completeness, condition~(iv) can also be expressed in terms
of the extension class. Choose a section
\(s:\Gamma\to E\), \(s(1)=1\), and write
\[
 s(\gamma)s(\eta)
 =\alpha(\gamma,\eta)s(\gamma\eta),\qquad
 \alpha(\gamma,\eta)\in A.
\]
If a map \(h:\Gamma\to V\) satisfies
\[
 \alpha(\gamma,\eta)
 =h(\gamma)+\rho(\gamma)h(\eta)-h(\gamma\eta),
\]
then \(b(a s(\gamma))=a+h(\gamma)\) is the required
crossed homomorphism. Conversely, restricting such a crossed homomorphism \(b\)
to the section gives this formula. Thus~(iv) says precisely
that the extension class becomes zero in \(H^2(\Gamma;V)\).
The central-element construction is an explicit way of obtaining
this vanishing without calculating a cohomology group.

\section{Mapping class group quotients}
\label{sec:groups}

\subsection{The Torelli group and its lower central series}

All mapping classes are represented by orientation-preserving
homeomorphisms. Denote \(H_g=H_1(\Sigma_g;\Z)\). The algebraic
intersection pairing \(\omega:H_g\times H_g\to\Z\) is
unimodular and alternating. Choose a symplectic basis
\(e_1,\ldots,e_g,f_1,\ldots,f_g\), so that
\(\omega(e_i,f_j)=\delta_{ij}\) and the pairings among the
\(e_i\)'s and among the \(f_i\)'s vanish. Denote the algebraic
symplectic group by \(\Sp_{2g}\).
The homological action is a surjective homomorphism
\[
 \sigma:\Mod(\Sigma_g)\longrightarrow\Sp_{2g}(\Z).
\]
Its kernel is the Torelli group \(\mathcal{T}_g\), see \cite{FarbMargalit} for this and
other basic facts about mapping class groups. 

For \(c\ge1\), we write
\[
 \mathcal{Q}_{g,c}=\Mod(\Sigma_g)/(\mathcal{T}_g)_{[c+1]},
 \qquad
 \mathcal{N}_{g,c}=\mathcal{T}_g/(\mathcal{T}_g)_{[c+1]}.
\]
We assume \(g\ge3\).

Since \((\mathcal{T}_g)_{[j]}\) is characteristic in \(\mathcal{T}_g\), it is
normal in \(\Mod(\Sigma_g)\). The quotient \(\mathcal{Q}_{g,c}\) is therefore
defined without lifting the symplectic action to \(\mathcal{T}_g\).
A group is
\emph{nilpotent of class at most \(c\)} when its
\((c+1)\)-st lower-central subgroup is trivial. In particular,
\(\mathcal{N}_{g,c}\) has class at most \(c\).

\begin{lemma}
\label{lem:nilab}
For \(g\ge3\) and \(c\ge1\), \(\mathcal{N}_{g,c}\) is finitely
generated and nilpotent. Moreover,
\[
 (\mathcal{N}_{g,c})_{\ab}\cong \mathcal{T}_g/[\mathcal{T}_g,\mathcal{T}_g],
 \qquad \mathcal{Q}_{g,c}/[\mathcal{N}_{g,c},\mathcal{N}_{g,c}]\cong \mathcal{Q}_{g,1}.
\]
These identifications respect the conjugation actions induced
from \(\Mod(\Sigma_g)\).
\end{lemma}
\begin{proof}
The Torelli group is finitely generated for \(g\ge3\), by
\cite{JohnsonI}. Hence \(\mathcal{N}_{g,c}\) is finitely generated.
For a normal subgroup \(K\subseteq I\), the commutator
subgroup of \(I/K\) is \([I,I]K/K\). Here
\(K=I_{[c+1]}\subseteq[I,I]\), with equality when
\(c=1\). Consequently
\[
 [\mathcal{N}_{g,c},\mathcal{N}_{g,c}]=(\mathcal{T}_g)_{[2]}/(\mathcal{T}_g)_{[c+1]}.
\]
Taking quotients gives the two isomorphisms. Both are induced by
quotient maps from \(\Mod(\Sigma_g)\), so they respect the conjugation action.
\end{proof}

\subsection{The integral Johnson module}

Let
\[
 A_g=\frac{\mathcal{T}_g/[\mathcal{T}_g,\mathcal{T}_g]}
 {\tor(\mathcal{T}_g/[\mathcal{T}_g,\mathcal{T}_g])}
\]
be the torsion-free quotient of \(\mathcal{T}_g/[\mathcal{T}_g,\mathcal{T}_g]\).
Conjugation on \(\mathcal{T}_g/[\mathcal{T}_g,\mathcal{T}_g]\) is trivial for elements of \(\mathcal{T}_g\),
so it factors through \(\Sp_{2g}(\Z)\). The same is true on
\(\tor(\mathcal{T}_g/[\mathcal{T}_g,\mathcal{T}_g])\) and \(A_g\).

Let \[\Omega=\sum_{i=1}^g e_i\wedge f_i\] and
\[j(v)=v\wedge\Omega.\] This is an equivariant map
\(j:H_g\to\Lambda^3 H_g\).


The following is the torsion-free part of Johnson's theorem
\cite{JohnsonIII}, see also \cite[p.~398 and Theorem~3.1]{DimcaPapadima}.
\begin{theorem}[Johnson's abelianization theorem]
\label{thm:johnsoninput}
For \(g\ge3\), the Johnson homomorphism gives a surjective
\(\Mod(\Sigma_g)\)-equivariant homomorphism
\[
 \tau_J:\mathcal{T}_g\longrightarrow\Lambda^3 H_g/j(H_g)
\]
and induces an equivariant isomorphism
\[
 A_g\cong\Lambda^3 H_g/j(H_g).
\]
\end{theorem}

The distinction between \(\mathcal{T}_g/[\mathcal{T}_g,\mathcal{T}_g]\) and \(A_g\) will be used below:
the affine realization is defined after quotienting by \(\tor(\mathcal{T}_g/[\mathcal{T}_g,\mathcal{T}_g])\),
and property~\(\rt\) then passes back through this finite kernel.

\begin{lemma}\label{lem:lattice}
For \(g\ge2\), \(j\) is injective with primitive image.
Consequently \(A_g\) is a free abelian group of rank
\[
 d_g=\binom{2g}{3}-2g \qquad(g\ge3).
\]
In particular, \(d_3=14\), \(d_4=48\), and \(d_5=110\).
\end{lemma}
\begin{proof}
For a basis vector \(v=e_i\) or \(f_i\), the nonzero terms
of \(j(v)\) are \(v\wedge e_k\wedge f_k\), \(k\ne i\),
each with coefficient \(1\) or \(-1\) in the ordered exterior
basis. Such a triple has a unique full symplectic pair, and its
remaining singleton determines \(v\). Thus the supports of the
\(2g\) vectors \(j(v)\) are disjoint. Choose one term for
each \(v\), possible since \(g\ge2\), and take its signed
coefficient. These coefficient functionals define an integral map
\(p:\Lambda^3 H_g\to H_g\) with \(pj=1\). Hence \(j(H_g)\)
is an integral direct summand. The quotient is free and its rank
is the difference of the ranks.
\end{proof}

The splitting in the proof need not be symplectic-equivariant.
For the real representation we will use an elementary weight
calculation. We first specify the
meaning of weights in this argument.

\paragraph{Weights for a diagonal one-parameter subgroup.}
Assume \(g\ge3\). Write \(H_{g,\R}=H_g\otimes_\Z\R\), and extend \(j\) linearly
to \(H_{g,\R}\). Johnson's identification gives
\[
 V_g:=A_g\otimes_\Z\R
 \cong \Lambda^3 H_{g,\R}/j(H_{g,\R}).
\]
For \(t\in\R\), consider the linear transformation
\[
 a_t=
 \begin{pmatrix}
 e^t I_g&0\\
 0&e^{-t}I_g
 \end{pmatrix},
 \qquad
 a_t e_i=e^t e_i,\quad a_t f_i=e^{-t}f_i.
\]
It preserves the symplectic pairing, since
\(\omega(a_t e_i,a_t f_j)=e^t e^{-t}\delta_{ij}=\delta_{ij}\),
and the other pairings remain zero. Thus
\(a_t\in\Sp_{2g}(\R)\), and \[a_{t+s}=a_ta_s,\] 
providing a one-parameter subgroup that we need for our argument.
We use the same notation for the induced operators on exterior
powers and on their invariant quotients.

For such a representation \(V\) and \(k\in\R\), define
\[
 V(k)=\{v\in V:a_t v=e^{kt}v\text{ for every }t\in\R\}.
\]
If \(V(k)\ne0\), we call \(k\) a \emph{weight} and \(V(k)\)
its \emph{weight space}, always relative to this particular
one-parameter subgroup. Thus a weight space is a simultaneous
eigenspace for all the \(a_t\),  \(k\) is the exponent, whereas
the eigenvalue of \(a_t\) is \(e^{kt}\). For example, \(e_i\)
has weight \(1\), \(f_i\) has weight \(-1\), and weight zero
means being fixed by every \(a_t\).


\begin{lemma}
\label{lem:weights}
For \(g\ge3\), the induced action of \((a_t)_{t\in\R}\)
on \(V_g\) has the direct-sum decomposition
\[
 V_g=V_g(3)\oplus V_g(1)\oplus V_g(-1)\oplus V_g(-3).
\]
On \(V_g(k)\) it is multiplication by \(e^{kt}\), and
\(V_g(0)=0\). Moreover,
\(-I\in\Sp_{2g}(\Z)\) acts as \(-1\) on \(A_g\).
\end{lemma}
\begin{proof}
Set
\[
 E=\operatorname{span}_\R\{e_1,\ldots,e_g\},\qquad
 F=\operatorname{span}_\R\{f_1,\ldots,f_g\}.
\]
Then \(H_{g,\R}=E\oplus F\), and
\[
 W:=\Lambda^3 H_{g,\R}
 =\bigoplus_{r=0}^3
   \bigl(\Lambda^r E\wedge\Lambda^{3-r}F\bigr).
\]
The summand indexed by \(r\) is spanned by exterior triples
with \(r\) factors from \(E\) and \(3-r\) factors from \(F\).
The action on an exterior product is obtained by acting on
each factor. Consequently, every vector \(w\) in this summand
satisfies
\[
 a_t w=e^{rt}e^{-(3-r)t}w=e^{(2r-3)t}w.
\]
In other words, weights add under exterior products. The four
possibilities \(r=3,2,1,0\) give weights \(3,1,-1,-3\),
respectively.

We now check explicitly what happens upon taking the quotient
by \(j(H_{g,\R})\). Since
\[
 a_t\Omega
 =\sum_{i=1}^g(e^te_i)\wedge(e^{-t}f_i)
 =\Omega,
\]
the bivector \(\Omega\) has weight zero. Therefore
\[
 a_t j(v)=a_t(v\wedge\Omega)
         =(a_t v)\wedge\Omega=j(a_t v).
\]
It follows that \(j(E)\subset W(1)\) and
\(j(F)\subset W(-1)\). Moreover,
\[
 j(H_{g,\R})=j(E)\oplus j(F),
\]
because \(j\) is injective and \(H_{g,\R}=E\oplus F\).
Thus this subspace is already the direct sum of its intersections
with the weight spaces of \(W\). Taking the quotient gives the
following natural identifications:
\begin{center}
\begin{tabular}{ccl}
\toprule
Weight \(k\) & \(W(k)\) & \(V_g(k)\), naturally isomorphic to\\
\midrule
\(3\) & \(\Lambda^3E\) & \(\Lambda^3E\)\\
\(1\) & \(\Lambda^2E\wedge F\) &
  \((\Lambda^2E\wedge F)/j(E)\)\\
\(-1\) & \(E\wedge\Lambda^2F\) &
  \((E\wedge\Lambda^2F)/j(F)\)\\
\(-3\) & \(\Lambda^3F\) & \(\Lambda^3F\)\\
\bottomrule
\end{tabular}
\end{center}
The quotient therefore inherits a direct-sum decomposition into
these weight spaces and no weight-zero space appears.
Finally, \(-I\) acts by \((-1)^3=-1\) on \(\Lambda^3H_g\),
and the same holds on its quotient \(A_g\).
\end{proof}

For example, when \(g=3\), these four weight spaces have
dimensions \(1,6,6,1\). Indeed, the two pure exterior cubes
have dimension \(\binom33=1\), while each mixed summand
has dimension \(\binom32\cdot3=9\), from which the
three-dimensional subspace \(j(E)\) or \(j(F)\) is removed.
Their dimensions add to \(\dim V_3=14\).

The absence of weight zero is the point needed later.
If \(v\in V_g(k)\), then for any norm on \(V_g\),
\[
 \|a_t v\|=e^{kt}\|v\|.
\]
Hence a vector of positive weight is contracted to zero as
\(t\to-\infty\), and a vector of negative weight is contracted
as \(t\to+\infty\). Every vector is a sum of these weight
components, so \(V_g\) is spanned by vectors that can be
contracted by elements of \(\Sp_{2g}(\R)\). We do not require
one choice of direction in \(t\) to contract all components
simultaneously. This is precisely the spanning condition used
in \cref{prop:affineT,prop:arithmeticT}.

\subsection{The affine symplectic group}
\label{sec:affine}

We verify the hypotheses of the general strategy for the integral
Johnson module. The first issue is contraction; the second is
volume preservation, which is needed to pass to a lattice.

\begin{proposition}
\label{prop:arithmeticT}
For \(g\ge3\), the groups
\[
 V_g\rtimes\Sp_{2g}(\R)
 \quad\text{and}\quad A_g\rtimes\Sp_{2g}(\Z)
\]
have property~\(\rt\).
\end{proposition}
\begin{proof}
Let \(\rho:\Sp_{2g}(\R)\to\GL(V_g)\) denote the action.
The real group \(\Sp_{2g}(\R)\) has property~\(\rt\),
and \(\Sp_{2g}(\Z)\) is a lattice in it.
These are the classical higher-rank and arithmetic-lattice
theorems, see \cite[Chapters~1 and~3 and Appendix~B]{BHV}
and \cite{Raghunathan}.

By \cref{lem:weights}, \(V_g\) is the direct sum of its weight
spaces with weights \(3,1,-1,-3\). A vector of positive weight
\(k\) is contracted by \(a_{-n}\), since
\[
 \rho(a_{-n})v=e^{-kn}v\longrightarrow0.
\]
A vector of negative weight is contracted by \(a_n\). Thus
\cref{prop:affineT} proves property~\(\rt\) for
\(V_g\rtimes\Sp_{2g}(\R)\).

We next verify that volume is preserved, so that the lattice
criterion in \cref{prop:affinelattice} applies.
We verify this using determinants. 
A symplectic transformation \(s\) of \(H_{g,\R}\)
preserves the nonzero top-degree form \(\omega^{\wedge g}\).
Consequently
\[
 \det(s|_{H_{g,\R}})=1.
\]
For an invertible operator \(u\) on an \(n\)-dimensional vector
space,
\[
 \det(\Lambda^3u)=(\det u)^{\binom{n-1}{2}}.
\]
Indeed, after triangularizing over \(\C\), the diagonal entries
of \(\Lambda^3u\) are the products of triples of distinct
diagonal entries of \(u\). Each entry occurs in exactly
\(\binom{n-1}{2}\) such products. Applying this with \(n=2g\)
gives \(\det(\Lambda^3s)=1\).

The map \(j:H_{g,\R}\to\Lambda^3H_{g,\R}\) is injective and
equivariant, so the action on its image also has determinant
one. In a basis beginning with a basis of this invariant
subspace, the action on \(\Lambda^3H_{g,\R}\) is block
triangular. Its determinant is the product of the determinants
on the subspace and on the quotient. Since that quotient is
\(V_g\), we obtain
\[
 \det\rho(s)
 =\frac{\det(\Lambda^3s)}
        {\det((\Lambda^3s)|_{j(H_{g,\R})})}
 =1.
\]
Thus any Lebesgue measure \(\lambda\) on \(V_g\) satisfies
\[
 \lambda(\rho(s)C)
 =|\det\rho(s)|\,\lambda(C)=\lambda(C)
\]
for every Borel set \(C\subset V_g\). This equality is the required condition, needed in \cref{prop:affinelattice}.

The group \(\Sp_{2g}(\R)\) is unimodular. This can also be seen
by a determinant calculation: its adjoint representation on
the symplectic Lie algebra identifies with the action on
\(\operatorname{Sym}^2H_{g,\R}\), and
\[
 \det(\operatorname{Sym}^2s)=(\det s)^{2g+1}=1.
\]
Right translation by \(s\) multiplies a left-invariant volume
form by \(\det\Ad(s)^{-1}\); hence a left Haar measure
\(\nu\) on \(\Sp_{2g}(\R)\) is also right invariant.

For completeness, the finite-volume domain of
\cref{prop:affinelattice} has a concrete form here. Choose a
half-open parallelepiped \(P\) for \(A_g\) in \(V_g\), normalize
Lebesgue measure by \(\lambda(P)=1\), and choose a Borel
fundamental domain \(\mathcal D\) for the right action of
\(\Sp_{2g}(\Z)\) on \(\Sp_{2g}(\R)\), with
\(0<\nu(\mathcal D)<\infty\).
Then
\[
 \mathcal F=\{(\rho(s)p,s):p\in P,\ s\in\mathcal D\}
\]
is a fundamental domain for the right action of
\(A_g\rtimes\Sp_{2g}(\Z)\) on
\(V_g\rtimes\Sp_{2g}(\R)\). Each fiber
\(\rho(s)P\) has volume one, so Tonelli's theorem gives
\[
 (\lambda\times\nu)(\mathcal F)
 =\int_{\mathcal D}\lambda(\rho(s)P)\,d\nu(s)
 =\nu(\mathcal D)<\infty.
\]
The restriction of product Haar measure to \(\mathcal F\)
defines the invariant quotient measure by the standard
quotient-measure construction used in
\cref{prop:affinelattice}. Thus \(A_g\rtimes\Sp_{2g}(\Z)\) is a
lattice in the real affine group and inherits property~\(\rt\).
\end{proof}

\subsection{The first Torelli quotient}

Let
\[
 D^2\mathcal{T}_g=\{i\in\mathcal{T}_g:i^m\in[\mathcal{T}_g,\mathcal{T}_g]
                              \text{ for some }m\ge1\}.
\]
It is the inverse image of \(\tor(\mathcal{T}_g/[\mathcal{T}_g,\mathcal{T}_g])\) under \(\mathcal{T}_g\to\mathcal{T}_g/[\mathcal{T}_g,\mathcal{T}_g]\),
and hence is characteristic in \(\mathcal{T}_g\). Define
\[
 E_g=\overline{\mathcal{Q}}_{g,1}=\Mod(\Sigma_g)/D^2\mathcal{T}_g=\mathcal{Q}_{g,1}/\tor(\mathcal{T}_g/[\mathcal{T}_g,\mathcal{T}_g]).
\]
There is an exact sequence
\begin{equation}\label{eq:torsionfreeabelianextension}
 1\longrightarrow A_g\longrightarrow E_g
 \xrightarrow{p}\Sp_{2g}(\Z)\longrightarrow1.
\end{equation}


We now consider an affine realization of \(E_g\) proved by Hain and Matsumoto.
The next statement is the closed-surface case of
\cite[Lemma~6.3]{HainMatsumoto}, and also the specialization of
\cref{prop:centralaffine} to \(z=-I\). We spell out the construction
in this case to identify the translation lattice and its index.

\begin{proposition}[The Hain--Matsumoto affine realization]
\label{prop:affinequotient}
There is an injective homomorphism
\[
 \phi:E_g\longrightarrow A_g\rtimes\Sp_{2g}(\Z)
\]
which induces the identity on the symplectic quotient and
multiplication by \(2\) on \(A_g\). Its image has index
\(2^{d_g}\). In particular, \(E_g\) is linear over \(\Z\)
and has property~\(\rt\).
\end{proposition}
\begin{proof}
Write \(E=E_g\) and \(A=A_g\).
Choose a lift \(u\in E\) of the central element \(z=-I\)
of \(\Sp_{2g}(\Z)\). Its action on \(A\) is multiplication by
\(-1\), by \cref{lem:weights}. For \(x\in E\) let
\[
 b(x)=[x,u]=xux^{-1}u^{-1}.
\]
Since \(z\) is central in \(\Sp_{2g}(\Z)\), this commutator lies
in \(A\). We write the law of \(A\) additively. The identity
\[
 [xy,u]=x[y,u]x^{-1}[x,u]
\]
and commutativity in \(A\) imply
\[
 b(xy)=b(x)+p(x)b(y).
\]
For \(a\in A\), the action of \(u\) gives
\(b(a)=a-za=2a\). Therefore
\(\phi(x)=(b(x),p(x))\) is a homomorphism to the stated
semidirect product. If it kills \(x\), then \(p(x)=1\),
so \(x=a\in A\), and \(2a=0\). Since \(A\) is torsion-free,
we obtain \(a=0\), which proves injectivity.

The subgroup \(\phi(E_g)\) surjects onto \(\Sp_{2g}(\Z)\) under the second-coordinate projection, and its intersection
with the translation subgroup is \(2A\). Every coset of the
image in \(A\rtimes\Sp_{2g}(\Z)\) therefore has a representative
in \(A\); two such representatives determine the same coset
exactly when they differ by \(2A\). The index is
\([A:2A]=2^{d_g}\).

Choose a basis of \(A\), and
let \[r_g:\Sp_{2g}(\Z)\to\GL_{d_g}(\Z)\] 
be the action. The map
\[
 (a,S)\longmapsto
 \begin{pmatrix}
  r_g(S)&a&0\\
  0&1&0\\
  0&0&S
 \end{pmatrix}
 \quad\in\GL_{d_g+1+2g}(\Z)
\]
is an injective homomorphism. The lower block recovers \(S\),
and the translation column recovers \(a\), without requiring
faithfulness of \(r_g\). Finally, \cref{prop:arithmeticT}
and finite-index permanence give property~\(\rt\) of \(E\).
\end{proof}

\begin{remark}
An extension by the module \(A\) is described by a class
\(e\in H^2(\Sp_{2g}(\Z);A)\). Choose a set-theoretic
section \(s:\Sp_{2g}(\Z)\to E\) with \(s(1)=1\), and write
\[
 s(\gamma)s(\eta)=\alpha(\gamma,\eta)s(\gamma\eta),
 \qquad \alpha(\gamma,\eta)\in A.
\]
Associativity implies that \(\alpha\) is a \(2\)-cocycle;
changing \(s\) changes it by a coboundary. Put
\(v(\gamma)=b(s(\gamma))\). Applying the crossed-homomorphism
identity and \(b|_A=2\) to the displayed product gives
\[
 2\alpha(\gamma,\eta)
 =v(\gamma)+\gamma v(\eta)-v(\gamma\eta).
\]
Thus \(2e=0\). Pushing out along \(A\xrightarrow{\times2}A\)
makes the extension split, yet it does not show that \(e=0\), and
does not give a section of \eqref{eq:torsionfreeabelianextension}.
The embedding of \cref{prop:affinequotient} corresponds to this
enlargement of the normal lattice.
\end{remark}

\begin{corollary}\label{cor:firststage}
The group \(\mathcal{Q}_{g,1}=\Mod(\Sigma_g)/[\mathcal{T}_g,\mathcal{T}_g]\) has property~\(\rt\)
for \(g\ge3\).
\end{corollary}
\begin{proof}
The kernel \(\tor(\mathcal{T}_g/[\mathcal{T}_g,\mathcal{T}_g])\) of \(\mathcal{Q}_{g,1}\to E_g\) is finite.
Since \(E_g\) has property~\(\rt\) by \cref{prop:affinequotient},
permanence under finite normal kernels gives the conclusion.
\end{proof}


\subsection{Higher Torelli quotients and proof of Theorem B}

We can now prove the mapping class group application.

\begin{proof}[Proof of \cref{thm:main}]
By \cref{lem:nilab}, \(\mathcal{N}_{g,c}\) is nilpotent and normal
in \(\mathcal{Q}_{g,c}\), and
\[
 \mathcal{Q}_{g,c}/[\mathcal{N}_{g,c},\mathcal{N}_{g,c}]
 \cong \mathcal{Q}_{g,1}.
\]
The latter has property~\(\rt\) by \cref{cor:firststage}.
The nilpotent permanence theorem, \cref{thm:nilpotentpermanence},
therefore gives property~\(\rt\) of \(\mathcal{Q}_{g,c}\).
Finite generation is obvious.
By \cref{lem:nilab,lem:lattice}, the abelianization of \(\mathcal{N}_{g,c}\)
surjects onto the nonzero free abelian group \(A_g\), of rank
\(d_g>0\). Thus \(\mathcal{N}_{g,c}\), and hence \(\mathcal{Q}_{g,c}\), is infinite.
\end{proof}

Note that
\(\mathcal{N}_{g,c}\) itself does not have property~\(\rt\), since it
surjects onto \(\Z^{d_g}\). 

\begin{corollary}[Quotients with nilpotent Torelli image]
\label{cor:nilpotentimage}
Let \(f:\Mod(\Sigma_g)\to P\), with \(g\ge3\), be a
surjective homomorphism. If \(f(\mathcal{T}_g)\) is nilpotent of finite
class, then \(P\) has property~\(\rt\).
\end{corollary}
\begin{proof}
If the class is at most \(c\ge1\), then
\(f\bigl((\mathcal{T}_g)_{[c+1]}\bigr)
=\bigl(f(\mathcal{T}_g)\bigr)_{[c+1]}=1\).
Thus \(f\) factors through \(\mathcal{Q}_{g,c}\), and quotient
permanence applies. The trivial-image case factors through the
symplectic quotient and is covered as well.
\end{proof}

\begin{remark}
For \(g\ge3\) and \(c\ge1\), the torsion subgroup of
\(\mathcal{N}_{g,c}\) is finite and characteristic
\cite[Chapter~II]{Raghunathan}, hence normal in \(\mathcal{Q}_{g,c}\).
Thus \(\mathcal{Q}_{g,c}/\tor(\mathcal{N}_{g,c})\) also has
property~\(\rt\), with torsion-free nilpotent Torelli image.
This is the quotient obtained by replacing \((\mathcal{T}_g)_{[c+1]}\)
by its rational lower-central analogue: the elements of \(\mathcal{T}_g\)
having a positive power in \((\mathcal{T}_g)_{[c+1]}\).
\end{remark}

There are natural surjections
\[
 \cdots\longrightarrow \mathcal{Q}_{g,3}\longrightarrow \mathcal{Q}_{g,2}
 \longrightarrow \mathcal{Q}_{g,1}\longrightarrow\Sp_{2g}(\Z).
\]
Increasing \(c\) retains more of the Torelli group. The theorem
applies separately to each finite nilpotent stage.


\subsection{An obstruction for genus 2}
\label{sec:genustwo}

The hypothesis \(g\ge3\) in \cref{thm:main} cannot be removed.
In genus \(g=2\), every quotient in the corresponding family has a
finite-index subgroup with infinite abelianization. We use the same
notation
\[
 \mathcal{Q}_{2,c}=\Mod(\Sigma_2)/(\mathcal{T}_2)_{[c+1]},\qquad
 \mathcal{N}_{2,c}=\mathcal{T}_2/(\mathcal{T}_2)_{[c+1]},\qquad c\ge1.
\]
The obstruction comes from a theorem of McCarthy concerning the
first cohomology of level subgroups.

\begin{proposition}\label{prop:genustwo}
For every integer \(c\ge1\), the group \(\mathcal{Q}_{2,c}\) has a
finite-index subgroup admitting a surjective homomorphism onto
\(\Z\). In particular, \(\mathcal{Q}_{2,c}\) does not have
property~\(\rt\).
\end{proposition}
\begin{proof}
Let
\[
 \Gamma=\Mod(\Sigma_2)[2]
 =\ker\bigl(\Mod(\Sigma_2)\longrightarrow\Sp_4(\Z/2\Z)\bigr)
\]
be the level-two mapping class group. It has finite index in
\(\Mod(\Sigma_2)\) and contains \(\mathcal{T}_2\), since the Torelli group acts
trivially on integral first homology and hence on homology modulo
two. McCarthy's theorem \cite[Theorem~B; Theorem~2.1]{McCarthy}
gives
\[
 H^1(\Gamma;\Z)\ne0.
\]
Here the coefficients are trivial, so
\(H^1(\Gamma;\Z)=\Hom(\Gamma,\Z)\). Choose a nonzero
homomorphism \(\chi:\Gamma\to\Z\). Its image is an infinite
cyclic subgroup of \(\Z\). Identifying that image with \(\Z\),
we may assume that \(\chi\) is surjective.

Put \(K_c=(\mathcal{T}_2)_{[c+1]}\). Since \(c\ge1\), we have
\[
 K_c\subseteq[\mathcal{T}_2,\mathcal{T}_2]
     \subseteq[\Gamma,\Gamma]
     \subseteq\ker\chi.
\]
Thus \(\chi\) induces a surjection
\[
 \bar\chi:\Gamma/K_c\longrightarrow\Z.
\]
The subgroup \(K_c\) is normal in \(\Mod(\Sigma_2)\) and contained in
\(\Gamma\). Consequently \(\Gamma/K_c\) is a subgroup of
\(\mathcal{Q}_{2,c}\), and
\[
 [\mathcal{Q}_{2,c}:\Gamma/K_c]=[\Mod(\Sigma_2):\Gamma]<\infty.
\]
If \(\mathcal{Q}_{2,c}\) had property~\(\rt\), then its finite-index
subgroup \(\Gamma/K_c\), and hence its quotient \(\Z\), would
also have property~\(\rt\), which is impossible.
\end{proof}

The argument shows why passing to any of these nilpotent Torelli
quotients leaves the obstruction intact. Since
\(K_c\subseteq[\Gamma,\Gamma]\), the quotient map induces an
isomorphism
\[
 (\Gamma/K_c)_{\ab}\cong\Gamma_{\ab}.
\]
Thus the character supplied by McCarthy survives at every stage.
No finite-generation assumption on \(\mathcal{T}_2\) is needed.

\begin{corollary}\label{cor:genustworelative}
For every \(c\ge1\), the pair \((\mathcal{Q}_{2,c},\mathcal{N}_{2,c})\) does not
have relative property~\(\rt\).
\end{corollary}
\begin{proof}
The subgroup \(\mathcal{N}_{2,c}\) is normal in \(\mathcal{Q}_{2,c}\), and
\[
 \mathcal{Q}_{2,c}/\mathcal{N}_{2,c}\cong\Sp_4(\Z).
\]
The latter group has property~\(\rt\). If the pair had relative
property~\(\rt\), then \cref{lem:normalextension} would give
property~\(\rt\) for \(\mathcal{Q}_{2,c}\), contrary to
\cref{prop:genustwo}.
\end{proof}

The symplectic quotient \(\Mod(\Sigma_2)/\mathcal{T}_2\cong\Sp_4(\Z)\) therefore
has property~\(\rt\), whereas none of the quotients
\(\mathcal{Q}_{2,c}\), \(c\ge1\), does. This establishes the necessity
of the genus assumption for the property~\(\rt\) theorem.

\section{Automorphisms of free groups}
\label{sec:freeautomorphisms}

We now apply the general criterion to the homological representation
of a free-group automorphism group. This application is independent
of the mapping class group calculation: the integral module, the
contracting diagonal subgroup and the arithmetic ambient group are
different. The same argument nevertheless proves property~\(\rt\)
for every quotient in which the kernel of the homological action
has nilpotent image. We consider general rank \(n\ge3\) in the
preliminary calculations and specialize to \(n=3\) in the proof
of \cref{thm:autquotients}.

\subsection{The homological representation and its kernel}

Let \(F_n=\langle x_1,\ldots,x_n\rangle\) be the free group of
rank \(n\), and let
\[
 H_n=H_1(F_n;\Z)=F_n/(F_n)_{[2]}\cong\Z^n,
 \qquad H_n^*=\Hom(H_n,\Z).
\]
Write \(e_i=[x_i]\) for the standard basis of \(H_n\), and
\(e_i^*\) for its dual basis. Abelianization induces a surjection
\[
 q:\Aut(F_n)\longrightarrow\GL_n(\Z).
\]
Indeed, the automorphisms which permute the free generators, invert
one generator, or replace \(x_i\) by \(x_ix_j\), for \(i\ne j\),
induce the corresponding integral elementary operations. These
operations generate \(\GL_n(\Z)\). Define
\[
 \operatorname{IA}_n=\ker q,
 \qquad
 \operatorname{SAut}(F_n)=q^{-1}(\SL_n(\Z)).
\]
Thus \(\operatorname{IA}_n\) consists of automorphisms acting
trivially on first homology, and \(\operatorname{SAut}(F_n)\)
has index two in \(\Aut(F_n)\). The latter is a determinant
condition.

For \(c\ge1\), define
\[
 \mathcal{P}_{n,c}
   =\Aut(F_n)/(\operatorname{IA}_n)_{[c+1]}.
\]
The subgroup \((\operatorname{IA}_n)_{[c+1]}\) is characteristic
in \(\operatorname{IA}_n\), hence normal in \(\Aut(F_n)\).
Consequently there is an exact sequence
\begin{equation}\label{eq:autnilsequence}
 1\longrightarrow
 \operatorname{IA}_n/(\operatorname{IA}_n)_{[c+1]}
 \longrightarrow\mathcal{P}_{n,c}
 \longrightarrow\GL_n(\Z)\longrightarrow1.
\end{equation}
No splitting of this sequence is assumed.

We prove \cref{thm:autquotients} below, beginning with the integral
abelianization of \(\operatorname{IA}_n\).

\subsection{The first Johnson homomorphism}

The commutator convention \([x,y]=xyx^{-1}y^{-1}\) gives the
standard identification
\[
 (F_n)_{[2]}/(F_n)_{[3]}
   \cong\bigwedge\nolimits^2H_n,
 \qquad [x,y]\longmapsto[x]\wedge[y],
\]
where \([x]\) and \([y]\) denote the images of \(x\) and \(y\) in the abelianization and the commutator on the left denotes its class modulo
\((F_n)_{[3]}\). This identification follows from the
commutator identities modulo third-order commutators: the
commutator is bilinear and alternating on abelianization, and
the classes \([x_j,x_k]\), \(j<k\), form a free abelian basis.
See \cite[Sections~3 and~4.1]{SatohSurvey}.

If \(\phi\in\operatorname{IA}_n\), then
\(\phi(x)x^{-1}\in(F_n)_{[2]}\) for every \(x\in F_n\).
Define
\begin{equation}\label{eq:autjohnsondefinition}
 \tau(\phi)([x])
   =[\phi(x)x^{-1}]
   \in(F_n)_{[2]}/(F_n)_{[3]}
   =\bigwedge\nolimits^2H_n.
\end{equation}
We check the conditions needed to regard this as an integral
homomorphism.

First, for \(x,y\in F_n\),
\[
 \phi(xy)(xy)^{-1}
 =\bigl(\phi(x)x^{-1}\bigr)
    x\bigl(\phi(y)y^{-1}\bigr)x^{-1}.
\]
Conjugation by \(x\) is trivial on
\((F_n)_{[2]}/(F_n)_{[3]}\). Thus the class of this expression
is the sum of the classes for \(x\) and \(y\). It follows that
\(x\mapsto[\phi(x)x^{-1}]\) is a homomorphism from \(F_n\)
to an abelian group, and therefore factors through \(H_n\).
This also proves that \eqref{eq:autjohnsondefinition} is
independent of the representative of \([x]\).

Second, every element of \(\operatorname{IA}_n\) acts trivially
on \((F_n)_{[2]}/(F_n)_{[3]}\): under the exterior-square
identification its action is the exterior square of the identity
on \(H_n\). For \(\phi,\psi\in\operatorname{IA}_n\), the
identity
\[
 (\phi\psi)(x)x^{-1}
   =\phi\bigl(\psi(x)x^{-1}\bigr)
      \phi(x)x^{-1}
\]
therefore gives \[\tau(\phi\psi)=\tau(\phi)+\tau(\psi).\]
We obtain a homomorphism
\[
 \tau:\operatorname{IA}_n\longrightarrow W_n,
 \qquad
 W_n=\Hom\left(H_n,\bigwedge\nolimits^2H_n\right)
     =H_n^*\otimes\bigwedge\nolimits^2H_n.
\]
Finally, for \(\theta\in\Aut(F_n)\), direct substitution in
\eqref{eq:autjohnsondefinition} gives
\begin{equation}\label{eq:autjohnsonequivariance}
 \tau(\theta\phi\theta^{-1})
    =\left(\bigwedge\nolimits^2q(\theta)\right)\,
        \tau(\phi)\,q(\theta)^{-1}.
\end{equation}
Thus \(\tau\) is equivariant for the conjugation action and
the natural \(\GL_n(\Z)\)-action on \(W_n\).

\begin{proposition}[Integral abelianization of
\(\operatorname{IA}_n\)]\label{prop:iaabelianization}
For \(n\ge3\), the homomorphism \(\tau\) induces an isomorphism
of \(\GL_n(\Z)\)-modules
\[
 \operatorname{IA}_n/[\operatorname{IA}_n,\operatorname{IA}_n]
       \cong W_n.
\]
In particular, its kernel is
\([\operatorname{IA}_n,\operatorname{IA}_n]\), and the
abelianization has no torsion. Its rank is
\[
 d_n=n\binom n2=\frac{n^2(n-1)}2.
\]
\end{proposition}

This is the standard computation of the first Johnson
homomorphism and the abelianization of \(\operatorname{IA}_n\),
see \cite[Sections~4.2--4.3.1, especially equation~(4.3)]{SatohSurvey}
for the statement and its historical attributions. We recall how
the integral generators enter the computation. Magnus's theorem
\cite[Section~6]{Magnus1935} (see also
\cite[Theorem~4.1]{SatohSurvey}) states that
\(\operatorname{IA}_n\) is generated by
\[
 \begin{aligned}
 C_{ij}:&\quad x_i\longmapsto x_j^{-1}x_ix_j
       &&(i\ne j),\\
 M_{ijk}:&\quad x_i\longmapsto x_i[x_j,x_k]
       &&(i,j,k\text{ distinct},\ j<k),
 \end{aligned}
\]
with all other free generators fixed. In particular,
\(\operatorname{IA}_n\) is finitely generated. Their images
are
\[
 \tau(C_{ij})=e_i^*\otimes(e_i\wedge e_j),
 \qquad
 \tau(M_{ijk})=e_i^*\otimes(e_j\wedge e_k).
\]
Up to signs, these are exactly the basis vectors
\(e_i^*\otimes(e_j\wedge e_k)\), \(j<k\), of \(W_n\).
This explains surjectivity and the absence of a finite-index
ambiguity in the integral image. Equivariance is already given
by \eqref{eq:autjohnsonequivariance}; the assertion about the
kernel is the abelianization statement in
\cref{prop:iaabelianization}.

\subsection{The affine realization}

By \cref{prop:iaabelianization}, the first quotient fits into the short exact sequence
\[
 1\longrightarrow W_n\longrightarrow\mathcal{P}_{n,1}
   \longrightarrow\GL_n(\Z)\longrightarrow1.
\]
The central element \(-I\in\GL_n(\Z)\) acts as \(-1\) on
\(H_n^*\) and as \(+1\) on \(\bigwedge^2H_n\). Consequently
it acts as \(-1\) on \(W_n\). This remains true when \(n\)
is odd: in that case \(-I\) is outside \(\SL_n(\Z)\), but it
is still central in the full quotient \(\GL_n(\Z)\).

\begin{proposition}\label{prop:autaffinerealization}
There is an injective homomorphism
\[
 \mathcal{P}_{n,1}\longrightarrow W_n\rtimes\GL_n(\Z)
\]
whose image has index \(2^{d_n}\).
\end{proposition}
\begin{proof}
Apply \cref{prop:centralaffine} to the central element
\(-I\), for which \(1-(-I)=2I\) is injective on \(W_n\)
and has cokernel of order \(2^{d_n}\).

More explicitly, choose a lift \(t\in\mathcal{P}_{n,1}\) of
\(-I\), and define \(b(u)=[u,t]\). Centrality in the quotient
implies that \(b(u)\in W_n\), even though \(t\) need not be
central in \(\mathcal{P}_{n,1}\). The commutator identity
\[
 [uv,t]=u[v,t]u^{-1}[u,t]
\]
and the fact that \(W_n\) is abelian give
\[b(uv)=b(u)+q(u)b(v).\] 
Thus
\(u\mapsto(b(u),q(u))\) is a homomorphism to the semidirect
product. For \(w\in W_n\), written additively,
\[
 b(w)=w-twt^{-1}=2w.
\]
The kernel of the homomorphism is therefore trivial. Its image
projects onto \(\GL_n(\Z)\), and its intersection with the
translation subgroup is exactly \(2W_n\). Every coset of the
image has a representative in the translation subgroup, so the
index is \([W_n:2W_n]=2^{d_n}\).
\end{proof}


\subsection{Contracting weights and the arithmetic lattice}

Let
\[
 V_n=W_n\otimes\R,
 \qquad
 S_n=\{s\in\GL_n(\R):|\det s|=1\}.
\]
The representation \(\rho:S_n\to\GL(V_n)\) is given, in the
Hom description, by
\[
 \rho(s)f=\left(\bigwedge\nolimits^2s\right)f s^{-1}.
\]
The identity component of \(S_n\) is \(\SL_n(\R)\), of index
two. Hence \(S_n\) has property~\(\rt\) for \(n\ge3\).
It is also unimodular. Indeed, its modular homomorphism is
trivial on the semisimple group \(\SL_n(\R)\), and the
resulting homomorphism from the quotient of order two to
\(\R_{>0}\) must be trivial.

\begin{lemma}\label{lem:autweights}
For every \(n\ge3\), there is a diagonal one-parameter subgroup
\((a_t)_{t\in\R}\) of \(\SL_n(\R)\) such that every basis
vector \(e_i^*\otimes(e_j\wedge e_k)\), \(j<k\), of \(V_n\)
has nonzero weight. In particular, each such vector is
contracted either by \(\rho(a_t)\) as \(t\to+\infty\), or
by \(\rho(a_{-t})\) as \(t\to+\infty\).
\end{lemma}
\begin{proof}
Choose
\[
 \alpha_i=3^{i-1}\quad(1\le i<n),
 \qquad
 \alpha_n=-\sum_{i=1}^{n-1}3^{i-1},
 \qquad
 a_t=\operatorname{diag}
       (e^{\alpha_1t},\ldots,e^{\alpha_nt}).
\]
The exponents sum to zero, so \(\det a_t=1\). On the dual
basis vector \(e_i^*\), the multiplier is
\(e^{-\alpha_it}\), while on \(e_j\wedge e_k\) it is
\(e^{(\alpha_j+\alpha_k)t}\). Thus the weight in question is
\[
 \beta_{ijk}=-\alpha_i+\alpha_j+\alpha_k,
 \qquad j<k.
\]
We verify explicitly that it is nonzero.

If \(i=j\) or \(i=k\), cancellation leaves one of the
nonzero numbers \(\alpha_1,\ldots,\alpha_n\). Otherwise
\(i,j,k\) are distinct. If all three are less than \(n\),
equality \(\beta_{ijk}=0\) would express one power of three
as a sum of two distinct powers of three, which is impossible
by uniqueness of the base-three expansion. If \(i=n\), all
three summands \(-\alpha_n,\alpha_j,\alpha_k\) are positive.
Finally, if \(n\) is one of \(j,k\), let \(\ell\) denote the
other one. Then
\[
 \beta_{ijk}
 =-\alpha_i+\alpha_\ell
       -\sum_{r=1}^{n-1}\alpha_r<0,
\]
because \(i\ne\ell\) and all the numbers in this last sum
are positive. These cases exhaust the possibilities.

On a vector of weight \(\beta\), the action is multiplication
by \(e^{\beta t}\). If \(\beta<0\), take \(t\to+\infty\).
If \(\beta>0\), take \(t\to-\infty\). In both cases the
multiplier tends to zero.
\end{proof}

For example, when \(n=3\), the chosen exponents are
\((1,3,-4)\). If \(i,j,k\) are distinct, the trace-zero
condition gives \(\beta_{ijk}=-2\alpha_i\), which is
nonzero. If an index cancels, the weight is one of
\(1,3,-4\). Thus the rank-three case needed for
\cref{thm:autquotients} is covered by the same argument.

\begin{proposition}\label{prop:autaffineT}
For \(n\ge3\), the groups \[V_n\rtimes S_n \text{\ \ \ \  and\ \ \ \ \ }
W_n\rtimes\GL_n(\Z)\] have property~\(\rt\).
\end{proposition}
\begin{proof}
The vectors contracted in \cref{lem:autweights} form a basis
of \(V_n\). Applying \cref{prop:affineT} with the subgroup
\(S_n\) therefore proves property~\(\rt\) for
\(V_n\rtimes S_n\). In terms of unitary representations,
property~\(\rt\) of \(S_n\) first supplies an
\(S_n\)-invariant vector. The contraction argument forces
that vector to be invariant under each translation along a
basis vector, hence under all of \(V_n\).

To pass to the integral group, we check volume preservation.
Write \(r_n=\binom n2\), so \(\dim V_n=nr_n\). The
determinant identities for an exterior square and a tensor
product give
\[
 \begin{aligned}
 \det(\rho(s))
 &=\det(s^{-T})^{r_n}
       \det\left(\bigwedge\nolimits^2s\right)^{n}\\
 &=(\det s)^{-r_n}(\det s)^{n(n-1)}
   =(\det s)^{r_n}.
 \end{aligned}
\]
For completeness, in the exterior-square determinant each
eigenvalue of \(s\) occurs in exactly \(n-1\) pairs, so
\[\det\left(\bigwedge^2s\right)=(\det s)^{n-1}.\]
For operators
\(u\) and \(v\) on spaces of dimensions \(p\) and \(q\),
respectively, the eigenvalues of \(u\otimes v\) are their
pairwise products, yielding
\[\det(u\otimes v)=(\det u)^q(\det v)^p.\]
These identities can be established by triangularizing over
\(\C\), without assuming diagonalizability.

Since \(s\in S_n\), we conclude
\[|\det\rho(s)|=1.\] Thus \(\rho(s)\) preserves Lebesgue
measure on \(V_n\), including when it reverses orientation.
The group \(\GL_n(\Z)\) preserves the lattice \(W_n\), and
is a lattice in \(S_n\). The latter assertion follows from
the classical lattice \(\SL_n(\Z)\subseteq\SL_n(\R)\) by passing
to these finite extensions, see \cite[Appendix~B]{BHV}.


All hypotheses of \cref{prop:affinelattice} now hold. More
explicitly, normalize Lebesgue measure so that a fundamental
parallelepiped for \(W_n\) has volume one, and multiply it
by Haar measure on \(S_n\). This is Haar measure on the
affine group. If \(\mathcal D\) is a finite-volume right
fundamental domain for \(\GL_n(\Z)\) in \(S_n\), the affine
fundamental domain has fiber \(\rho(s)P\) over \(s\in
\mathcal D\), where \(P\) is the parallelepiped. Each fiber
has volume one, so the total volume is the Haar measure of
\(\mathcal D\), which is finite. Thus
\(W_n\rtimes\GL_n(\Z)\) is a lattice in the real affine
group and inherits property~\(\rt\).
\end{proof}

\subsection{Passing to higher nilpotent stages}

\begin{proof}[Proof of \cref{thm:autquotients}]
By \cref{prop:autaffinerealization,prop:autaffineT}, the group
\(\mathcal{P}_{3,1}\) is isomorphic to a finite-index
subgroup of a group with property~\(\rt\), and hence has
property~\(\rt\).

Fix \(c\ge1\) and let \[J=\operatorname{IA}_3/
(\operatorname{IA}_3)_{[c+1]}.\] Magnus's finite-generation
theorem \cite[Section~6]{Magnus1935} shows that \(J\) is finitely generated. Its
\((c+1)\)-st lower-central subgroup is trivial, and therefore
it is nilpotent of class at most \(c\). Furthermore,
\[
 [J,J]=(\operatorname{IA}_3)_{[2]}/
                (\operatorname{IA}_3)_{[c+1]},
 \qquad
 \mathcal{P}_{3,c}/[J,J]\cong\mathcal{P}_{3,1}.
\]
The nilpotent quotient theorem \cref{thm:quotienttower} gives
property~\(\rt\) for \(\mathcal{P}_{3,c}\). Equivalently,
all hypotheses of \cref{thm:generalcriterion} have been
verified with \(G=\Aut(F_3)\), \(K=\operatorname{IA}_3\),
\(A=W_3\), \(\Gamma=\GL_3(\Z)\), \(S=S_3\), and the
central element \(-I\).

The same commutator identity gives
\(J/[J,J]\cong W_3\), which is an infinite free abelian
group. Thus \(J\) is infinite. The group
\(\mathcal{P}_{3,c}\) is finitely generated because
\(\Aut(F_3)\) is generated by finitely many Nielsen
automorphisms, and is infinite also because it surjects onto
\(\GL_3(\Z)\).

Finally, suppose that \(f:\Aut(F_3)\to E\) is a surjective
homomorphism for which \(f(\operatorname{IA}_3)\) is nilpotent
of class at most \(c\). Then
\(f((\operatorname{IA}_3)_{[c+1]})=1\), so \(f\) factors
through \(\mathcal{P}_{3,c}\). Quotient permanence proves
property~\(\rt\) for \(E\).
\end{proof}

\subsection{The Andreadakis filtration and the examples of
Lubotzky--Pak}

There is a related, but generally different, family of quotients
obtained by allowing automorphisms to act on nilpotent quotients
of the free group itself. For \(r\ge1\), define
\[
 \mathcal A_n(r)=\ker\!\left(
 \Aut(F_n)\longrightarrow
 \Aut\bigl(F_n/(F_n)_{[r+1]}\bigr)\right).
\]
This is the Andreadakis filtration \cite{Andreadakis}. Its first term is
\(\operatorname{IA}_n\), and
\cref{prop:iaabelianization} gives
\(\mathcal A_n(2)=[\operatorname{IA}_n,\operatorname{IA}_n]\).
Its centrality means that
\[
 [\mathcal A_n(r),\mathcal A_n(s)]
      \subseteq\mathcal A_n(r+s),
 \qquad
 (\operatorname{IA}_n)_{[r]}\subseteq\mathcal A_n(r).
\]
The first inclusion is the standard commutator property of the
Andreadakis filtration; see \cite[Section~3.2]{SatohSurvey}.
The second follows by induction from the first, starting with
\(\mathcal A_n(1)=\operatorname{IA}_n\).

\begin{corollary}\label{cor:tamefreequotients}
For \(n\ge3\) and \(r\ge1\), the image
\[
 \im\!\left(\Aut(F_n)\longrightarrow
        \Aut\bigl(F_n/(F_n)_{[r+1]}\bigr)\right)
\]
has property~\(\rt\). The corresponding image of
\(\operatorname{SAut}(F_n)\) has property~\(\rt\) as well.
\end{corollary}
\begin{proof}
For \(r=1\), the two images are \(\GL_n(\Z)\) and
\(\SL_n(\Z)\). For \(r\ge2\), the inclusion
\((\operatorname{IA}_n)_{[r]}\subseteq\mathcal A_n(r)\)
gives a surjection
\[
 \mathcal{P}_{n,r-1}\longrightarrow
        \Aut(F_n)/\mathcal A_n(r).
\]
The group on the right is the image in the statement. For \(n=3\),
it has property~\(\rt\) by \cref{thm:autquotients}; for \(n\ge4\),
this follows from property~\(\rt\) of \(\Aut(F_n)\) and quotient
permanence \cite{KalubaNowakOzawa,KalubaKielakNowak,NitscheAutFour}.
The special image is its subgroup of index two, detected by determinant
on first homology, and consequently also has property~\(\rt\).
\end{proof}

The special-image assertion of \cref{cor:tamefreequotients}
was proved by Lubotzky and Pak
\cite[Theorem~3.8 and Section~4]{LubotzkyPak}. Their proof
realizes that image as a lattice in a semidirect product of
a real unipotent group and \(\SL_n(\R)\), and proves
property~\(\rt\) for this ambient group. Thus the use of
arithmetic affine groups in this setting has a direct
precedent. Conversely, their theorem at \(r=2\), together
with \(\mathcal A_n(2)=[\operatorname{IA}_n,
\operatorname{IA}_n]\) and \cref{thm:quotienttower}, also
proves \cref{thm:autquotients}. The explicit first-stage
argument above identifies all of the ingredients without
using the higher-stage theorem of \cite{LubotzkyPak}.

The image in \cref{cor:tamefreequotients} is called the
\emph{tame automorphism group} of the free nilpotent group.
 In general, an automorphism of
\(F_n/(F_n)_{[r+1]}\) need not lift to an automorphism of
\(F_n\). Lubotzky and Pak explicitly distinguish the image from
the full automorphism group and note that the image has
infinite index for \(n\ge2\) and \(r\ge4\)
\cite[Example~1.14]{LubotzkyPak}.


\section{Nonlinearity}
\label{sec:nonlinearity}

We now prove \cref{thm:nonlinearity}. The obstruction is a rational
central subgroup in the unipotent completion of \(\mathcal{T}_g\). By a theorem
of Hain, this subgroup lies in the kernel of the homomorphism to the
unipotent radical of the relative completion of \(\Mod(\Sigma_g)\). It survives,
however, in the unipotent completion of each lower-central-series quotient
\(\mathcal N_{g,c}\), for \(c\ge2\). We will show
that a finite-kernel complex representation of \(\mathcal{Q}_{g,c}\) would kill
infinitely many elements of \(\mathcal{N}_{g,c}\) corresponding to this central
subgroup.

There are two points to verify. First, we must arrange that the image
of the whole Torelli group is unipotent, without introducing an infinite
kernel. Second, we must pass from the central subgroup in the completion
to elements of the discrete group. The first uses property~\(\rt\)
and the structure of nilpotent algebraic groups. The second uses the
rationality of Hain's central subgroup and Malcev theory.

\subsection{Algebraic closures and finite abelianization}

Unless a field of definition is specified, algebraic groups in the
characteristic-zero argument are affine algebraic groups over \(\C\).
We denote them by bold letters. Recall that a linear algebraic group
is a Zariski-closed subgroup of some \(\GL_d\), and that its identity
component \(\mathbf H^\circ\) has finite index. A matrix is unipotent
if all its eigenvalues are equal to \(1\), and semisimple if it is
diagonalizable over the algebraic closure. A torus is an algebraic
group isomorphic over that field to a product of multiplicative groups.

The unipotent radical \(\Ru(\mathbf H^\circ)\) is the largest
connected normal unipotent subgroup of \(\mathbf H^\circ\). The
quotient is reductive. We use the facts that a quotient by a closed
normal algebraic subgroup is again affine, that the derived subgroup
of a connected linear algebraic group is closed, and that a connected
nilpotent algebraic group in characteristic zero is a product of a
torus and a unipotent group. See \cite{Borel} and, for the nilpotent
Jordan decomposition, \cite[Section~2.1 and Appendix~A]{Baues}.

A group \(\Gamma\) has \emph{FAb} if every finite-index subgroup of
\(\Gamma\) has finite abelianization \cite[Section~1]{AbertLubotzkyPyber}. A group is \emph{perfect} if it
equals its commutator subgroup. For a connected algebraic group we use
the algebraic derived subgroup in this definition.

A discrete group with property~\(\rt\) has FAb. Indeed, property
\(\rt\) passes to finite-index subgroups and their abelianizations,
and a discrete abelian group with property~\(\rt\) is finite
\cite[Corollary~1.7.2]{BHV}. In particular, \cref{thm:main} gives FAb
for \(\mathcal{Q}_{g,c}\). 

\begin{lemma}\label{lem:fabclosure}
Let \(\Gamma\) have FAb, let \(r:\Gamma\to\GL_d(\C)\) be a
homomorphism, and let \(\mathbf H\) be the Zariski closure of its
image. Then \(\mathbf H^\circ\) is perfect and has no nontrivial
central torus.
\end{lemma}
\begin{proof}
Set \(\Gamma_0=r^{-1}(\mathbf H^\circ)\). This is a finite-index
subgroup of \(\Gamma\), and \(r(\Gamma_0)\) is Zariski dense in
\(\mathbf H^\circ\): each component is both open and closed, so
intersection with a dense subgroup is dense in that component.

The image of \(\Gamma_0\) in
\[
 \mathbf H^\circ/[\mathbf H^\circ,\mathbf H^\circ]
\]
is finite, since \(\Gamma_0\) has finite abelianization. It is also
Zariski dense. A finite set is closed, and hence this connected
algebraic quotient is finite and therefore trivial.

Let \(\mathbf U=\Ru(\mathbf H^\circ)\). The reductive group
\(\mathbf H^\circ/\mathbf U\) is perfect, so it is semisimple and
has finite center. A central torus \(\mathbf T\) of
\(\mathbf H^\circ\) has connected image in that center. Its image
is therefore trivial, and \(\mathbf T\subseteq\mathbf U\). Since an
element which is both semisimple and unipotent is the identity, we
obtain \(\mathbf T=1\).
\end{proof}

The conclusion concerns central tori, not the whole center. In
particular, a perfect algebraic group may have a nontrivial unipotent
center.

\begin{lemma}\label{lem:unipotentization}
Under the hypotheses of \cref{lem:fabclosure}, let
\(N\subseteq\Gamma\) be a nilpotent normal subgroup and let \(\mathbf K\) be
the Zariski closure of \(r(N)\). There are a finite normal algebraic
subgroup \(\mathbf F\subseteq\mathbf H\) and a connected
unipotent normal algebraic subgroup \(\mathbf U\) such that
\[
 \mathbf K=\mathbf F\times\mathbf U.
\]
Consequently the homomorphism
\[
 \bar r:\Gamma\xrightarrow{r}\mathbf H(\C)
       \longrightarrow(\mathbf H/\mathbf F)(\C)
\]
maps all of \(N\) into a connected unipotent normal subgroup.
If \(r\) has finite kernel, then so does \(\bar r\).
\end{lemma}
\begin{proof}
The group \(r(\Gamma)\) normalizes \(\mathbf K\), and Zariski
density implies that \(\mathbf H\) does as well. Suppose that \(N\)
has nilpotency class at most \(c\). Every iterated commutator of
length \(c+1\) is trivial on \(r(N)\). The commutator word is a
regular map, and a product of Zariski-dense sets is Zariski dense.
Thus the same identity holds on \(\mathbf K\), so \(\mathbf K\)
is nilpotent.

Write \(\mathbf K^\circ=\mathbf T\times\mathbf U\), where
\(\mathbf T\) is a torus and \(\mathbf U\) is unipotent. The
torus consists of the semisimple elements of \(\mathbf K^\circ\)
and is preserved by algebraic automorphisms. In particular,
\(\mathbf H^\circ\) normalizes it. A normal torus in a connected
algebraic group is central: conjugation acts on its discrete character
lattice, and such an action of a connected group is trivial. By
\cref{lem:fabclosure}, \(\mathbf T=1\), and therefore
\(\mathbf K^\circ=\mathbf U\).

The characteristic-zero Levi decomposition gives
\(\mathbf K=\mathbf F\ltimes\mathbf U\), where \(\mathbf F\)
is finite because \(\mathbf K/\mathbf U\) is finite. For
\(f\in\mathbf F\), nilpotence implies that \(\Ad(f)\) acts
unipotently on \(\Lie\mathbf U\). The map \(u\mapsto[f,u]\)
fixes the identity and has derivative \(\Ad(f)-I\) there.
Since \(\mathbf K\) has nilpotency class at most \(c\), its
\(c\)-fold composition is constant. The chain rule therefore gives
\((\Ad(f)-I)^c=0\) on \(\Lie\mathbf U\).
On the other hand, \(f\) has finite order,
so \(\Ad(f)\) is semisimple in characteristic zero. Hence
\(\Ad(f)=1\). Exponential and logarithm identify \(\mathbf U\)
with its nilpotent Lie algebra, and it follows that \(f\) centralizes
\(\mathbf U\). Thus
\[
 \mathbf K=\mathbf F\times\mathbf U.
\]

By uniqueness of Jordan decomposition, \(\mathbf F\) is precisely
the set of semisimple elements of \(\mathbf K\). It is therefore
preserved by every algebraic automorphism of \(\mathbf K\), including
conjugation by \(\mathbf H\). In particular, \(\mathbf F\) is
normal in the whole algebraic group \(\mathbf H\); see also
\cite[Section~2.1 and Appendix~A]{Baues}.

The image of \(\mathbf U\) in \(\mathbf H/\mathbf F\) is
connected, unipotent and normal, and contains the image of \(N\).
Finally, \(\ker\bar r=r^{-1}(\mathbf F)\), so
\[
 |\ker\bar r|\le |\mathbf F|\,|\ker r|
\]
whenever \(\ker r\) is finite.
\end{proof}

It is important here to quotient by \(\mathbf F\) before passing to
the identity component. Passing to a component first could replace
\(N\) by a proper finite-index subgroup. The lemma puts the image of
all of \(N\) in the identity component. We may then take finite-index
subgroups in the symplectic quotient while maintaining the full Torelli
group, as required below.

\subsection{Malcev completion}

Let \(\Ga\) denote the additive algebraic group, so that
\(\Ga(\Q)=(\Q,+)\). A prounipotent group is an inverse limit of
unipotent algebraic groups. The rational unipotent completion
\(\Lambda^{\un}\) of a finitely generated group \(\Lambda\) is
equipped with a homomorphism
\(\Lambda\to\Lambda^{\un}(\Q)\) universal for homomorphisms
from \(\Lambda\) to rational unipotent algebraic groups. Its Lie
algebra is the rational Malcev Lie algebra of \(\Lambda\).
Extension of scalars to \(\C\) gives the corresponding universal
property for complex unipotent targets.

We recall the form of Malcev theory which we will use. See
\cite[Chapter~II, Theorems~2.12 and~2.18, and Remark~2.16]{Raghunathan}
and \cite[Section~4]{HainMatsumoto} for details.

\begin{theorem}[Malcev]\label{thm:malcev}
Let \(N\) be a finitely generated nilpotent group. Its torsion
elements form a finite characteristic subgroup \(\tor(N)\). There
is a unipotent algebraic group \(\mathbf N\) over \(\Q\) and a
canonical homomorphism
\[
 j:N\longrightarrow\mathbf N(\Q),\qquad \ker j=\tor(N),
\]
such that \(j(N)\) is a lattice in \(\mathbf N(\R)\): it is
discrete and has compact quotient. The group \(\mathbf N\) is the
unipotent completion of \(N\).

If \(\mathbf V\subseteq\mathbf N\) is an algebraic subgroup defined over \(\Q\), then
\(j(N)\cap\mathbf V(\R)\) is a lattice in \(\mathbf V(\R)\).
In particular, if \(\mathbf V\cong\Ga\), this intersection is
infinite cyclic.

If \(N=\Lambda/\Lambda_{[c+1]}\), then
\[
 \Lie\mathbf N\cong
 \mathfrak m(\Lambda)/\mathfrak m(\Lambda)_{[c+1]},
\]
where \(\mathfrak m(\Lambda)\) is the rational Malcev Lie algebra
of \(\Lambda\).

\end{theorem}

On the Lie-algebra side we use the analogous notation for the lower
central series, taking closed ideals in the pronilpotent case.
The last identification also follows from the universal property.
A homomorphism from \(\Lambda\) into a unipotent group of class at
most \(c\) kills \(\Lambda_{[c+1]}\), and the corresponding
quotient of its Malcev Lie algebra is
\(\mathfrak m(\Lambda)/\mathfrak m(\Lambda)_{[c+1]}\).

\begin{lemma}\label{lem:rationalkernel}
With the notation of \cref{thm:malcev}, let
\(\mathbf V\subseteq\mathbf N\) be a rational subgroup isomorphic to
\(\Ga\), and let \(b:N\to\mathbf W(\C)\) be a homomorphism
into a complex unipotent algebraic group. Denote its algebraic
extension by \(b^{\un}:\mathbf N_{\C}\to\mathbf W\).
If \(b^{\un}\) is trivial on \(\mathbf V_{\C}\), then
\(\ker b\) is infinite.
\end{lemma}
\begin{proof}
By \cref{thm:malcev}, choose a nonidentity element
\(j(n)\in j(N)\cap\mathbf V(\R)\). It has infinite order because
\(\mathbf V(\R)\cong(\R,+)\). We have
\[
 b(n)=b^{\un}(j(n))=1.
\]
The powers of \(n\) are distinct, since their images under \(j\)
are distinct, and all lie in \(\ker b\).
\end{proof}

Rationality is essential to this argument. An irrational line in a
real vector space may meet an integral lattice only at zero. Zariski
density alone is therefore not enough to imply the conclusion of the lemma.

\subsection{Relative completion}

Write the homological representation as \(\sigma:\Mod(\Sigma_g)\to\Sp_{2g}(\Q)\). Its image
\(\Sp_{2g}(\Z)\) is Zariski dense. The relative unipotent completion
of \(\Mod(\Sigma_g)\) with respect to \(\sigma\) is an extension
\[
 1\longrightarrow\mathcal U_g\longrightarrow\mathcal G_g
 \longrightarrow\Sp_{2g,\Q}\longrightarrow1,
\]
where \(\mathcal U_g\) is prounipotent, together with a lift
\(\widetilde\sigma:\Mod(\Sigma_g)\to\mathcal G_g(\Q)\) of \(\sigma\).
It has the following universal property. If
\[
 1\longrightarrow\mathbf W\longrightarrow\mathbf E
 \longrightarrow\Sp_{2g,\Q}\longrightarrow1
\]
is an algebraic extension with \(\mathbf W\) unipotent, every
homomorphism \(\Mod(\Sigma_g)\to\mathbf E(\Q)\) lifting \(\sigma\) factors
through a morphism \(\mathcal G_g\to\mathbf E\). For \(g\ge3\),
this relative completion commutes with extension of scalars from
\(\Q\) to \(\C\), by \cite[Theorem~3.1]{Hain}. Thus the
same universal property holds for complex targets. This
completion retains the prescribed symplectic representation and allows
unipotent information above it. It also differs from the ordinary unipotent
completion of \(\Mod(\Sigma_g)\). See \cite[Section~3]{Hain} and
\cite[Section~4]{HainMatsumoto}. 

Let \(\mathscr{P}_g\) denote the rational unipotent completion
of \(\mathcal{T}_g\), and let \(\mathfrak t_g\) be its Lie algebra.
Write \(\mathscr{P}_{g,\C}\) for its extension of scalars to \(\C\).
The image of \(\mathcal{T}_g\) in \(\mathcal G_g\) lies
in \(\mathcal U_g\), since \(\sigma(\mathcal{T}_g)=1\). It induces a
homomorphism \(\theta:\mathscr{P}_g\to\mathcal U_g\). For
\(m\ge1\), set
\[
 \Sp_{2g}(\Z)[m]=\ker\bigl(\Sp_{2g}(\Z)
                  \to\Sp_{2g}(\Z/m\Z)\bigr),\qquad
 \Mod(\Sigma_g)[m]=\sigma^{-1}(\Sp_{2g}(\Z)[m]).
\]
In particular, \(\mathcal{T}_g\) is a subgroup of every \(\Mod(\Sigma_g)[m]\):
\[
 \mathcal{T}_g\subseteq\Mod(\Sigma_g)[m]\qquad(m\ge1).
\]

\begin{theorem}[Hain]\label{thm:hainkernel}
For \(g\ge3\), the following statements hold over \(\Q\).
\begin{enumerate}[label=(\roman*)]
\item Restriction of \(\widetilde\sigma\) to \(\Mod(\Sigma_g)[m]\) gives
      its completion relative to \(\sigma|_{\Mod(\Sigma_g)[m]}\), for every
      \(m\ge1\). Thus passing from \(\Mod(\Sigma_g)\) to any level subgroup
      \(\Mod(\Sigma_g)[m]\) does not change its relative completion
      with respect to the homological representation.
\item There is an exact central extension
      \[
       1\longrightarrow\mathbf Z_g\longrightarrow\mathscr{P}_g
       \xrightarrow{\theta}\mathcal U_g\longrightarrow1,
       \qquad \mathbf Z_g\cong\Ga.
      \]
\item If \(\mathfrak z_g=\Lie\mathbf Z_g\), then
      \[
       \mathfrak z_g\subset(\mathfrak t_g)_{[2]},
       \qquad
       \mathfrak z_g\cap(\mathfrak t_g)_{[3]}=0.
      \]
      Consequently this rational central line survives in
      \(\mathfrak t_g/(\mathfrak t_g)_{[c+1]}\) for every
      \(c\ge2\).
\end{enumerate}
\end{theorem}

These statements are Proposition~3.3 (p.~604), Theorem~3.4 (p.~605),
and Theorem~4.10 and the ensuing discussion (pp.~611--612) of
\cite{Hain}. For (iii), Hain identifies the central line with
\(\Q(1)\), the one-dimensional Tate Hodge structure of weight
\(-2\), and identifies the weight filtration with the lower central
filtration. We use only the resulting inclusions in (iii). In
particular, the argument does not require a quadratic presentation
of the Torelli Lie algebra in a larger genus range.

\begin{corollary}\label{cor:centralline}
Let \(g\ge3\) and \(c\ge2\), and let \(\mathbf N_c\) be the
rational unipotent completion of \(\mathcal{N}_{g,c}\). The natural map
\(p_c:\mathscr{P}_g\to\mathbf N_c\) is injective on
\(\mathbf Z_g\). Its image \(\mathbf Z_c\) is a rational central
subgroup of \(\mathbf N_c\) isomorphic to \(\Ga\).
\end{corollary}
\begin{proof}
By \cref{thm:malcev},
\[
 \Lie\mathbf N_c=
 \mathfrak t_g/(\mathfrak t_g)_{[c+1]}.
\]
Since \(c+1\ge3\), part (iii) of \cref{thm:hainkernel} implies
that the quotient map is injective on \(\mathfrak z_g\).
Exponential identifies one-dimensional unipotent groups with their
Lie algebras in characteristic zero, so \(p_c\) is injective on
\(\mathbf Z_g\). Its image is defined over \(\Q\) and is central
because \(p_c\) is surjective.
\end{proof}

Thus the rational unipotent completion of the finite-stage Torelli quotient contains a rational central
subgroup which maps trivially into every unipotent extension of the
prescribed symplectic representation. To apply this observation to an
arbitrary complex representation, we need arithmetic superrigidity.

\subsection{Arithmetic extension}
We record the following statement will allow to complete the nonlinearity argument.
\begin{theorem}\label{thm:arithmeticextension}
Let \(g\ge3\).
\begin{enumerate}[label=(\roman*)]
\item Every finite-index subgroup of \(\Sp_{2g}(\Z)\) contains
      a principal congruence subgroup \(\Sp_{2g}(\Z)[m]\).
\item Let \(\Lambda\subseteq\Sp_{2g}(\Z)\) be a subgroup of finite index, let
      \(\mathbf R\) be a complex linear algebraic group, and let
      \(a:\Lambda\to\mathbf R(\C)\) be a homomorphism. There are
      a finite-index subgroup \(\Lambda'\subseteq\Lambda\) and an
      algebraic homomorphism
      \(\alpha:\Sp_{2g,\C}\to\mathbf R\) such that
      \(a(s)=\alpha(s)\) for all \(s\in\Lambda'\).
\end{enumerate}
\end{theorem}

Part (i) is the congruence subgroup theorem of Bass, Milnor and Serre
\cite[Theorem~14.1]{BassMilnorSerre}. Part (ii) is the virtual
algebraic-extension consequence of arithmetic superrigidity for the
integral symplectic group; see \cite[Chapter~VIII, Theorem~B]{Margulis}
and \cite[Section~6.1, Remarks after Conjecture~6.1, item~2]{ChurchFarb}.

We explain how to obtain the stated form from the matrix formulation.
For a representation of a finite-index subgroup \(\Lambda\), induce
to \(\Sp_{2g}(\Z)\) and apply virtual algebraic extension. The
original summand is invariant under a finite-index symplectic subgroup,
and hence under the algebraic symplectic group by Zariski density.
Restriction to that summand gives the required extension. Next, embed
\(\mathbf R\) as a closed subgroup of \(\GL_d\). An algebraic
extension with values in \(\GL_d\) takes values in \(\mathbf R\),
since it does so on a Zariski-dense finite-index subgroup. This gives
(ii) for an arbitrary complex linear algebraic target. By (i), the
finite-index restrictions can all be taken at principal congruence
levels.

\subsection{Proof of \texorpdfstring{\cref{thm:nonlinearity}}{Theorem D}}

\begin{proof}[Proof of \cref{thm:nonlinearity}]
Fix \(g\ge3\) and \(c\ge2\), and let
\[
 q:\Mod(\Sigma_g)\to\mathcal{Q}_{g,c}
\]
be the natural surjection.
Assume that \(r:\mathcal{Q}_{g,c}\to\GL_d(\C)\) has finite kernel. The
homological representation \(\sigma:\Mod(\Sigma_g)\to\Sp_{2g}(\Z)\)
factors through \(\mathcal{Q}_{g,c}\), since \((\mathcal{T}_g)_{[c+1]}\subseteq\mathcal{T}_g\),
and induces the quotient homomorphism
\(\mathcal{Q}_{g,c}\to\mathcal{Q}_{g,c}/\mathcal{N}_{g,c}=\Sp_{2g}(\Z)\).

\emph{Step 1.} By \cref{thm:main}, \(\mathcal{Q}_{g,c}\) has property~\(\rt\)
and hence FAb. Let \(\mathbf H=\overline{r(\mathcal{Q}_{g,c})}^{\mathrm{Zar}}\).
Applying \cref{lem:fabclosure,lem:unipotentization} gives a finite
normal subgroup \(\mathbf F\subseteq\mathbf H\) such that
\[
 \bar r:\mathcal{Q}_{g,c}\longrightarrow\overline{\mathbf H}(\C),
 \qquad \overline{\mathbf H}=\mathbf H/\mathbf F,
\]
has finite kernel and maps \(\mathcal{N}_{g,c}\) into a connected unipotent normal
subgroup \(\overline{\mathbf U}\) of \(\overline{\mathbf H}\).

\emph{Step 2.} Put \(\mathbf H_1=\overline{\mathbf H}^{\circ}\).
Connectedness of \(\overline{\mathbf U}\) gives
\(\bar r(\mathcal{N}_{g,c})\subseteq\mathbf H_1\). Therefore the homomorphism
\[
 \mathcal{Q}_{g,c}\longrightarrow\overline{\mathbf H}/\mathbf H_1
\]
factors through \(\mathcal{Q}_{g,c}/\mathcal{N}_{g,c}=\Sp_{2g}(\Z)\) and has finite
image. By \cref{thm:arithmeticextension}(i), there is \(m_0\ge1\)
such that
\[
 \bar r(q(\Mod(\Sigma_g)[m_0]))\subseteq\mathbf H_1(\C).
\]
The subgroup \(\Mod(\Sigma_g)[m_0]\) still contains all of \(\mathcal{T}_g\).

Let \(\mathbf R=\mathbf H_1/\Ru(\mathbf H_1)\), and denote
the quotient map by \(p:\mathbf H_1\to\mathbf R\). Since
\(\overline{\mathbf U}\) is connected, unipotent and normal in
\(\mathbf H_1\), it lies in \(\Ru(\mathbf H_1)\). Thus
\(p\bar r q\) kills \(\mathcal{T}_g\) and defines a homomorphism
\[
 a:\Sp_{2g}(\Z)[m_0]\longrightarrow\mathbf R(\C),
 \qquad a(\sigma(f))=p(\bar r(q(f))).
\]
This is well defined because \(\ker\sigma=\mathcal{T}_g\).

\emph{Step 3.} By \cref{thm:arithmeticextension}(ii), followed by
(i), there are an integer \(m\) divisible by \(m_0\) and an
algebraic homomorphism \(\alpha:\Sp_{2g,\C}\to\mathbf R\)
such that
\[
 p(\bar r(q(f)))=\alpha(\sigma(f))\qquad(f\in \Mod(\Sigma_g)[m]).
\]
Consider the fiber product
\[
 \mathbf P=\{(s,h)\in\Sp_{2g,\C}\times\mathbf H_1:
                       \alpha(s)=p(h)\}.
\]
This is a linear algebraic group. Projection to the first factor is
surjective because \(p\) is surjective, and its kernel is
\(\{1\}\times\Ru(\mathbf H_1)\). Hence we obtain an algebraic
extension
\begin{equation}\label{eq:pullbackextension}
 1\longrightarrow\Ru(\mathbf H_1)\longrightarrow\mathbf P
 \longrightarrow\Sp_{2g,\C}\longrightarrow1.
\end{equation}
The map
\[
 \psi:\Mod(\Sigma_g)[m]\longrightarrow\mathbf P(\C),\qquad
 \psi(f)=(\sigma(f),\bar r(q(f)))
\]
is a homomorphism lifting the standard symplectic representation.
The fiber product puts the algebraic extension of the quotient
representation into the form required by relative completion: the
quotient is \(\Sp_{2g,\C}\) and the kernel is unipotent.

\emph{Step 4.} By \cref{thm:hainkernel}(i), the relative completion
of \(\Mod(\Sigma_g)[m]\) is \(\mathcal G_g\). By \cite[Theorem~3.1]{Hain},
its relative completion over \(\C\) is \(\mathcal G_{g,\C}\).
The universal property therefore gives a morphism
\[
 \Psi:\mathcal G_{g,\C}\longrightarrow\mathbf P,
\]
whose composition with the canonical map from \(\Mod(\Sigma_g)[m]\) is \(\psi\).
The restriction to \(\mathcal U_{g,\C}\) has image in
\(\Ru(\mathbf H_1)\), the kernel of \eqref{eq:pullbackextension}.
Consequently the homomorphism on \(\mathcal{T}_g\) factors as
\[
 \mathcal{T}_g\longrightarrow\mathscr{P}_g(\C)
 \xrightarrow{\theta}\mathcal U_g(\C)
 \longrightarrow\Ru(\mathbf H_1)(\C).
\]
It follows that the induced homomorphism on \(\mathscr{P}_{g,\C}\)
kills \(\mathbf Z_{g,\C}\).

\emph{Step 5.} The restriction \(b=\bar r|_{\mathcal{N}_{g,c}}\) has unipotent
image and therefore extends uniquely to an algebraic homomorphism
\[
 b^{\un}:\mathbf N_{c,\C}\longrightarrow\Ru(\mathbf H_1).
\]
Both \(b^{\un}p_c\) and the homomorphism on
\(\mathscr{P}_{g,\C}\) from Step~4 induce \(\bar r q\) on
\(\mathcal{T}_g\). Uniqueness in the universal property of unipotent
completion identifies them. Thus \(b^{\un}p_c\) kills
\(\mathbf Z_{g,\C}\). By \cref{cor:centralline}, the image
\(\mathbf Z_c=p_c(\mathbf Z_g)\) is a nontrivial rational additive
subgroup of \(\mathbf N_c\), and \(b^{\un}\) is trivial on
\(\mathbf Z_{c,\C}\). Lemma~\ref{lem:rationalkernel} now implies
that \(\ker b\) is infinite. This contradicts the finite kernel of
\(\bar r\) in Step~1.

Therefore \(\mathcal{Q}_{g,c}\) admits no complex representation with finite kernel,
and in particular no faithful complex representation.
\end{proof}

The contradiction uses more than the existence of a central Lie
algebra direction. Property~\(\rt\) and superrigidity force the
representation to kill Hain's central subgroup, while its rationality
forces the kernel to contain infinitely many elements of the discrete
nilpotent quotient. Both features are needed.

\subsection{Other fields and finite-index subgroups}

\begin{corollary}\label{cor:charzero}
For \(g\ge3\) and \(c\ge2\), the group \(\mathcal{Q}_{g,c}\) is not
linear over any field of characteristic zero, and no finite-index
subgroup of \(\mathcal{Q}_{g,c}\) is complex-linear.
\end{corollary}
\begin{proof}
The group \(\mathcal{Q}_{g,c}\) is finitely generated. Suppose that it embeds
in \(\GL_d(k)\), where \(k\) has characteristic zero. The entries
of the images of a finite generating set and their inverses generate
a finitely generated subfield \(k_0\subset k\). Such a field embeds
in \(\C\): map a transcendence basis to algebraically independent
complex numbers and extend over the remaining finite algebraic
extension. This gives a faithful complex representation, contrary
to \cref{thm:nonlinearity}.

If a finite-index subgroup \(\Lambda\subseteq \mathcal{Q}_{g,c}\) has a faithful
complex representation, induce it to \(\mathcal{Q}_{g,c}\). The induced
representation is finite-dimensional and faithful. Indeed, an element
outside \(\Lambda\) moves the identity-coset summand, while an
element of \(\Lambda\) acts on that summand by the given faithful
representation. This again contradicts \cref{thm:nonlinearity}.
\end{proof}

\begin{lemma}\label{lem:positivechar}
A finitely generated nilpotent group linear over a field of positive
characteristic is virtually abelian.
\end{lemma}
\begin{proof}
Extend the field to its algebraic closure. Let \(L\subseteq\GL_d\) be
the given subgroup and let \(\mathbf L\) be its Zariski closure. The
commutator-identity argument in \cref{lem:unipotentization} shows
that \(\mathbf L\) is nilpotent. Set
\(L_0=L\cap\mathbf L^\circ\); this subgroup has finite index
in \(L\). A connected nilpotent algebraic group is a product of a
central torus and a unipotent group, also in positive characteristic
\cite[Section~10]{Borel}. Thus \([L_0,L_0]\) consists of unipotent
matrices.

If the characteristic is \(p\) and \(p^e\ge d\), every unipotent
matrix \(u\in\GL_d\) satisfies \(u^{p^e}=1\). To see this,
write \(u=1+X\), where \(X^d=0\), and use
\((1+X)^{p^e}=1+X^{p^e}\). The subgroup \([L_0,L_0]\) is finitely
generated, since subgroups of finitely generated nilpotent groups are
finitely generated. It is therefore a finitely generated nilpotent
torsion group, and hence finite. These nilpotent-group facts follow
from a finite central series with finitely generated abelian factors.

Finally, a finitely generated group with finite commutator subgroup
has center of finite index. Each generator has finitely many conjugates,
all in its coset modulo the commutator subgroup, so its centralizer has
finite index. Intersecting the centralizers of a finite generating set
proves the assertion. Thus \(L_0\), and consequently \(L\), is
virtually abelian.
\end{proof}

\begin{corollary}\label{cor:allfields}
For \(g\ge3\) and \(c\ge2\), \(\mathcal{Q}_{g,c}\) is not linear over
any field.
\end{corollary}
\begin{proof}
By \cref{cor:charzero}, it remains to consider positive characteristic.
A faithful representation of \(\mathcal{Q}_{g,c}\) restricts faithfully to
the finitely generated nilpotent subgroup \(\mathcal{N}_{g,c}\), which would
then be virtually abelian by \cref{lem:positivechar}.

Finite-index subgroups of a finitely generated nilpotent group have
the same rational Malcev Lie algebra: their torsion-free images are
commensurable lattices in the same Malcev group. A virtually abelian
nilpotent group consequently has abelian Malcev Lie algebra. However,
\cref{thm:hainkernel}(iii) gives
\[
 0\ne(dp_c)(\mathfrak z_g)\subseteq
 \bigl(\mathfrak t_g/(\mathfrak t_g)_{[c+1]}\bigr)_{[2]}.
\]
Thus the Malcev Lie algebra of \(\mathcal{N}_{g,c}\) is not abelian, a
contradiction.
\end{proof}

\section{Final remarks}
\label{sec:finalremarks}

We conclude by briefly describing further applications of the same method.
In each case the main point is to identify an abelianized kernel
to which the affine argument applies, or to reduce to a quotient
already considered above.

\subsection{Free products of free abelian groups}

Let
\[
 A=\Z^{n_1}*\cdots *\Z^{n_r},\qquad r\ge3,\quad n_i\ge3,
\]
and let \(G=\operatorname{Out}(A)\). Denote by \(K\) the kernel
of the action of \(G\) on \(H_1(A;\Z)\). Then
\[
 G/K_{[c+1]}\text{ has property }\rt
 \qquad\text{for every }c\ge1.
\]
To see this, pass to the finite-index subgroup \(G_0\) preserving
the conjugacy class of each free factor. Its homological quotient
is \(\prod_i\GL_{n_i}(\Z)\), and automorphisms of the individual
factors give a section. The kernel \(K\) is generated by partial
conjugations: one conjugates a free factor by an element of another
and fixes the remaining factors. The standard generating-set
description is recalled in \cite[Section~2.5]{SaleRAAG}.
It follows that \(K/[K,K]\) is an equivariant quotient of
\[
 \bigoplus_{i=1}^r(\Z^{n_i})^{\oplus(r-1)},
\]
with the \(i\)-th arithmetic factor acting through its standard
representation on the corresponding summands.

The contraction argument and lattice passage in
\cref{sec:generalaffine} give property~\(\rt\) for the semidirect
product of this module with \(\prod_i\GL_{n_i}(\Z)\).
Its quotient \(G_0/[K,K]\) therefore has property~\(\rt\).
Finite-index permanence and \cref{thm:quotienttower} give the
assertion. The same reasoning applies to \(\Aut(A)\).

These are further examples in which the original group does not
have property~\(\rt\). Indeed, the defining graph is a disjoint
union of cliques and satisfies the largeness criterion of
\cite[Theorem~2(1b)]{GuirardelSale}. Thus \(G\) has a finite-index
subgroup mapping onto a nonabelian free group. 

\subsection{Quotients associated with the Johnson kernel}

Let
\[
 \mathcal J_g=\ker\!\left(\mathcal T_g
       \xrightarrow{\tau_J}\Lambda^3H_g/j(H_g)\right).
\]
For \(g\ge4\), the method also gives property~\(\rt\) for
\[
 \Mod(\Sigma_g)/(\mathcal J_g)_{[c+1]},\qquad c\ge1.
\]
The additional inputs are finite generation of \(\mathcal J_g\),
proved in \cite[Theorem~A]{ChurchErshovPutman}, and unipotence of
the action of \(\mathcal T_g/\mathcal J_g\) on
\(H_1(\mathcal J_g;\Q)\), proved in
\cite[Theorem~A]{DimcaHainPapadima}.

For the first stage, put
\[
 E=\Mod(\Sigma_g)/[\mathcal J_g,\mathcal J_g],
 \qquad N=\mathcal T_g/[\mathcal J_g,\mathcal J_g].
\]
The abelian normal subgroup \(\mathcal J_g/[\mathcal J_g,\mathcal J_g]\)
has finite torsion, and the quotient of \(N\) by this subgroup
is the free abelian group \(\mathcal T_g/\mathcal J_g\).
Unipotence means that some power of the augmentation ideal of
\(\Q[\mathcal T_g/\mathcal J_g]\) annihilates
\(H_1(\mathcal J_g;\Q)\). After removing the finite torsion,
the same identity holds on the integral lattice. Since the
operators \(t-1\) describe commutators with lifts of \(t\),
it follows that \(N\) is finite-by-nilpotent. Moreover,
\[
 E/[N,N]\cong\Mod(\Sigma_g)/[\mathcal T_g,\mathcal T_g].
\]
The first-stage result for the Torelli group and
\cref{rem:nilpotentextensions} give property~\(\rt\) for \(E\).
Now apply \cref{thm:quotienttower} with \(K=\mathcal J_g\).

The nonlinearity argument extends to these quotients, already
when \(c=1\). The image of \(\mathcal J_g\) in \(\mathscr P_g\)
lies in its second lower-central subgroup, so
\((\mathcal J_g)_{[c+1]}\) maps into degree at least \(2c+2\ge4\).
Consequently Hain's central line survives in the Malcev completion
of the Torelli image. This image is finite-by-nilpotent at every
stage. After removing the finite torsion of
\(\mathcal J_g/(\mathcal J_g)_{[c+1]}\), its rational
lower-central factors are quotients of tensor powers of
\(H_1(\mathcal J_g;\Q)\). Unipotence supplies invariant flags
on these factors, refining them to a central series for the
whole Torelli image, whose quotient by the Johnson image is
abelian. The proof of
\cref{thm:nonlinearity} applies. A finite-kernel complex
representation before this removal would induce one afterwards
by quotienting its algebraic image closure by the image of the
finite normal subgroup. The positive-characteristic argument
of \cref{lem:positivechar,cor:allfields} applies as well.
Thus these groups admit no complex representation with finite
kernel and are not linear over any field.

The interest here is the use of a different natural normal
subgroup, rather than an unrelated source of property~\(\rt\):
after removal of a finite kernel, each quotient is a quotient
of a sufficiently deep Torelli lower-central stage.
In genus three, \cite[Theorems~A and~E]{GaifullinJohnson} gives
finite generation of \((\mathcal J_3)_{\ab}\) and unipotence
after restriction to a finite-index subgroup of
\(\mathcal T_3/\mathcal J_3\). Extending the full-action argument
to this case requires a further step; the conclusion above
is stated only for \(g\ge4\).

\subsection{Marked points and boundary components}

Let \(\Sigma_{g,r}^{p}\) have \(p\) marked points and \(r\)
boundary components, all fixed pointwise by mapping classes.
Define \(\mathcal T_{g,r}^{p}\) as the kernel of the action on
the first homology of the capped closed surface. This convention
is important when there is more than one boundary component.
For \(g\ge3\), the same argument gives property~\(\rt\) for
\[
 \Mod(\Sigma_{g,r}^{p})/
             (\mathcal T_{g,r}^{p})_{[c+1]},\qquad c\ge1.
\]
Indeed, writing \(H_{g,\Q}=H_g\otimes_\Z\Q\), one has
\[
 H_1(\mathcal T_{g,r}^{p};\Q)
 \cong
 \bigl(\Lambda^3H_{g,\Q}/j(H_{g,\Q})\bigr)
       \oplus H_{g,\Q}^{\oplus(p+r)}.
\]
See \cite[Propositions~3.5 and~4.6]{Hain}.
Finite generation, finite torsion in the integral abelianization,
and algebraicity of its torsion-free quotient are recorded in
\cite[proof of Proposition~3.4]{HainMatsumoto}.
Every summand has the required contraction property, and
\(-I\) acts as \(-1\). Thus \cref{thm:strategy} applies.
Hain's central-line and relative-completion results hold for
these groups as well \cite[Theorems~3.4 and~4.10]{Hain}, so
the proof of \cref{thm:nonlinearity} and the other-field
arguments extend when \(c\ge2\).

There is an additional geometric feature. For \(r>0\) and
\(c\ge2\), the boundary twists generate a central subgroup
isomorphic to \(\Z^r\) in the displayed quotient.
With one boundary, its twist acts on \(F_{2g}\) by conjugation
by \(\omega=\prod_i[a_i,b_i]\). It has infinite order on
\(F_{2g}/(F_{2g})_{[4]}\): for a nonzero power, the displacement
of \(a_1\) has nonzero degree-three class
\(m[\sum_i[a_i,b_i],a_1]\) in the free Lie algebra.
On the other hand, the third lower-central subgroup of the
Torelli group acts trivially on this nilpotent quotient, by
the commutator property of the Johnson filtration.
Capping all but one boundary and forgetting the marked points
detects the boundary twists independently. Thus the construction
also gives nonlinear property-\(\rt\) groups with central
free abelian subgroups of arbitrarily large rank.

\subsection{Outer automorphisms of the free group of rank three}

Set
\[
 \operatorname{IO}_3=
 \ker\!\left(\operatorname{Out}(F_3)\longrightarrow\GL_3(\Z)\right).
\]
The natural map \(\Aut(F_3)\to\operatorname{Out}(F_3)\) maps
\(\operatorname{IA}_3\) onto \(\operatorname{IO}_3\), and maps
their lower-central terms onto one another. Hence it induces
a surjection
\[
 \mathcal P_{3,c}\longrightarrow
 \operatorname{Out}(F_3)/(\operatorname{IO}_3)_{[c+1]}.
\]
These quotients have property~\(\rt\) by
\cref{thm:autquotients}. They are finitely generated as quotients of
\(\Aut(F_3)\), and infinite because they map onto \(\GL_3(\Z)\).
This places the free-group application in the outer-automorphism
setting, which is closer to the closed-surface case. Again,
the original group \(\operatorname{Out}(F_3)\) does not have
property~\(\rt\); see \cite{McCool,GrunewaldLubotzky}.

\subsection{Expansion and dependence on nilpotency class}

Suppose that \(G\) is finitely generated and that
\(G/K_{[c+1]}\) has property~\(\rt\) for every \(c\ge1\).
Fix a finite symmetric generating set \(S\) of \(G\) and fix \(c\).
Every finite quotient of \(G\) in which the image of \(K\)
is nilpotent of class at most \(c\) factors through the same
Kazhdan group \(G/K_{[c+1]}\).
Consequently the lazy random walks on these quotients, defined
by the images of \(S\) with multiplicities retained, have a
common positive spectral gap away from constant functions.
Thus their Cayley graphs form an expander family whenever
their orders tend to infinity; see \cite[Chapter~6]{BHV}.

The dependence of this gap on \(c\) is a separate question.
For free nilpotent groups, product-replacement applications
were established by \cite{LubotzkyPak}, and explicit
depth-dependent Kazhdan bounds for tame automorphism groups
appear in \cite[Theorem~3.14]{HadadExtensions}.
Obtaining comparable bounds for the larger quotients considered
here, or bounds independent of \(c\), would strengthen the
qualitative permanence argument. Property~\(\rt\) at each
stage does not itself give a uniform bound. Conversely,
failure of property~\(\rt\) for the original group does not
force the quotient gaps to tend to zero: uniform expansion
of a family of quotients is weaker than property~\(\rt\)
of the original group.

\subsection*{Acknowledgements}
The author was supported by the National Science Center Grant Maestro-13
UMO-2021/42/A/ST1/00306. 


The author used GPT-Sol 5.6 and GPT-Astra 6 during the development of this work to explore proof strategies and examples, examine mathematical arguments, and assist with drafting and revising the manuscript. 
The author assumes responsibility for all content, including the verification of the proofs and the final presentation.


\phantomsection
\addcontentsline{toc}{section}{References}
\begin{thebibliography}{CWMS18}
\raggedright

\bibitem[ALP03]{AbertLubotzkyPyber}
M.~Ab\'ert, A.~Lubotzky and L.~Pyber,
\emph{Bounded generation and linear groups},
Internat. J. Algebra Comput. \textbf{13} (2003), no.~4, 401--413.

\bibitem[A65]{Andreadakis}
S.~Andreadakis, \emph{On the automorphisms of free groups and free
nilpotent groups}, Proc. London Math. Soc. (3) \textbf{15} (1965), 239--268.

\bibitem[BMS67]{BassMilnorSerre}
H. Bass, J. Milnor, and J.-P. Serre,
\emph{Solution of the congruence subgroup problem for
\(\mathrm{SL}_n\) (\(n\ge3\)) and
\(\mathrm{Sp}_{2n}\) (\(n\ge2\))},
Publ. Math. Inst. Hautes \'Etudes Sci. \textbf{33} (1967), 59--137.

\bibitem[B04]{Baues}
O. Baues, \emph{Infra-solvmanifolds and rigidity of subgroups in
solvable linear algebraic groups}, Topology \textbf{43} (2004),
no.~4, 903--924.

\bibitem[BHV08]{BHV}
B. Bekka, P. de la Harpe, and A. Valette,
\emph{Kazhdan's Property (T)}, New Mathematical Monographs 11,
Cambridge University Press, 2008.

\bibitem[B91]{Borel}
A. Borel, \emph{Linear Algebraic Groups}, second edition,
Graduate Texts in Mathematics 126, Springer-Verlag, New York, 1991.

\bibitem[CWMS18]{CWMS}
I. Chatterji, D. Witte Morris, and R. Shah,
\emph{Relative Property (T) for Nilpotent Subgroups},
Documenta Math. \textbf{23} (2018), 353--382.

\bibitem[CEP22]{ChurchErshovPutman}
T. Church, M. Ershov, and A. Putman,
\emph{On finite generation of the Johnson filtrations},
J. Eur. Math. Soc. \textbf{24} (2022), no.~8, 2875--2914.


\bibitem[CF13]{ChurchFarb}
T. Church and B. Farb, \emph{Representation theory and homological
stability}, Adv. Math. \textbf{245} (2013), 250--314.

\bibitem[DGO17]{DGO}
F. Dahmani, V. Guirardel, and D. Osin,
\emph{Hyperbolically embedded subgroups and rotating families
in groups acting on hyperbolic spaces},
Mem. Amer. Math. Soc. \textbf{245} (2017), no.~1156, v+152 pp.


\bibitem[DHP14]{DimcaHainPapadima}
A. Dimca, R. Hain, and S. Papadima,
\emph{The abelianization of the Johnson kernel},
J. Eur. Math. Soc. \textbf{16} (2014), no.~4, 805--822.


\bibitem[DP13]{DimcaPapadima}
A. Dimca and S. Papadima,
\emph{Arithmetic group symmetry and finiteness properties of Torelli
groups}, Ann. of Math. (2) \textbf{177} (2013), 395--423.


\bibitem[FM12]{FarbMargalit}
B. Farb and D. Margalit,
\emph{A Primer on Mapping Class Groups}, Princeton Mathematical
Series 49, Princeton University Press, 2012.

\bibitem[G25]{GaifullinJohnson}
A. A. Gaifullin,
\emph{Finite generation of abelianizations of the genus 3 Johnson
kernel and the commutator subgroup of the Torelli group for
\(\operatorname{Out}(F_3)\)},
arXiv:2507.20710, revised 2026.

\bibitem[GL09]{GrunewaldLubotzky}
F.~Grunewald and A.~Lubotzky,
\emph{Linear representations of the automorphism group of a
free group},
Geom. Funct. Anal. \textbf{18} (2009), no.~5, 1564--1608.

\bibitem[GS18]{GuirardelSale}
V. Guirardel and A. Sale,
\emph{Vastness properties of automorphism groups of RAAGs},
J. Topol. \textbf{11} (2018), no.~1, 30--64.

\bibitem[H10]{HadadExtensions}
U. Hadad,
\emph{Kazhdan constants of group extensions},
Internat. J. Algebra Comput. \textbf{20} (2010), no.~5, 671--688.

\bibitem[H97]{Hain}
R. Hain, \emph{Infinitesimal presentations of the Torelli groups},
J. Amer. Math. Soc. \textbf{10} (1997), no.~3, 597--651.


\bibitem[HM09]{HainMatsumoto}
R. Hain and M. Matsumoto, \emph{Relative pro-\(\ell\)
completions of mapping class groups}, J. Algebra \textbf{321}
(2009), no.~11, 3335--3374.


\bibitem[H16]{Hull}
M. Hull, \emph{Small cancellation in acylindrically hyperbolic
groups}, Groups Geom. Dyn. \textbf{10} (2016), no.~4, 1077--1119.


\bibitem[J80]{Johnson1980}
D.~Johnson, \emph{An abelian quotient of the mapping class group
\(\mathcal I_g\)}, Math. Ann. \textbf{249} (1980), no.~3, 225--242.


\bibitem[J83]{JohnsonI}
D. Johnson, \emph{The structure of the Torelli group. I. A finite set
of generators for \(\mathcal I\)}, Ann. of Math. (2)
\textbf{118} (1983), 423--442.

\bibitem[J85]{JohnsonIII}
D. Johnson, \emph{The structure of the Torelli group. III.
The abelianization of \(\mathcal I\)}, Topology \textbf{24}
(1985), 127--144.

\bibitem[KKN21]{KalubaKielakNowak}
M.~Kaluba, D.~Kielak and P.~W.~Nowak,
\emph{On property \((T)\) for \(\Aut(F_n)\) and
\(\SL_n(\Z)\)},
Ann. of Math. (2) \textbf{193} (2021), no.~2, 539--562.

\bibitem[KNO19]{KalubaNowakOzawa}
M.~Kaluba, P.~W.~Nowak and N.~Ozawa,
\emph{\(\Aut(F_5)\) has property \((T)\)},
Math. Ann. \textbf{375} (2019), 1169--1191.

\bibitem[LP01]{LubotzkyPak}
A.~Lubotzky and I.~Pak,
\emph{The product replacement algorithm and Kazhdan's property
\((T)\)},
J. Amer. Math. Soc. \textbf{14} (2001), no.~2, 347--363.

\bibitem[M35]{Magnus1935}
W.~Magnus, \emph{\"Uber \(n\)-dimensionale Gittertransformationen},
Acta Math. \textbf{64} (1935), 353--367.

\bibitem[M91]{Margulis}
G. A. Margulis, \emph{Discrete Subgroups of Semisimple Lie Groups},
Ergebnisse der Mathematik und ihrer Grenzgebiete (3), vol.~17,
Springer-Verlag, Berlin, 1991.

\bibitem[M01]{McCarthy}
J. D. McCarthy, \emph{On the first cohomology group of cofinite
subgroups in surface mapping class groups}, Topology
\textbf{40} (2001), no.~2, 401--418.


\bibitem[M89]{McCool}
J.~McCool, \emph{A faithful polynomial representation of
\(\operatorname{Out} F_3\)},
Math. Proc. Cambridge Philos. Soc. \textbf{106} (1989), no.~2, 207--213.

\bibitem[N13]{Nica}
B. Nica, \emph{Linear groups---Malcev's theorem and Selberg's lemma},
arXiv:1306.2385, 2013.

\bibitem[N20]{NitscheAutFour}
M.~Nitsche,
\emph{Computer proofs for Property \((T)\), and SDP duality},
arXiv:2009.05134, revised 2022.

\bibitem[R72]{Raghunathan}
M. S. Raghunathan, \emph{Discrete Subgroups of Lie Groups},
Ergebnisse der Mathematik und ihrer Grenzgebiete, vol.~68,
Springer-Verlag, Berlin--Heidelberg--New York, 1972.


\bibitem[S24]{SaleRAAG}
A. Sale,
\emph{On virtual indicability and property (T) for outer
automorphism groups of RAAGs},
Groups Geom. Dyn. \textbf{18} (2024), no.~1, 147--190.


\bibitem[S16]{SatohSurvey}
T.~Satoh,
\emph{A survey of the Johnson homomorphisms of the automorphism
groups of free groups and related topics},
in \emph{Handbook of Teichm\"uller Theory}, Vol.~V (2016),
167--209.


\end{thebibliography}

\bigskip
\noindent

\noindent\textsc{Piotr W. Nowak}\\ \textsc{Institute of Mathematics of the Polish Academy of Sciences, Warsaw, Poland.}\\
\textit{Email:} \texttt{pnowak@impan.pl}\\[0.5em]


\end{document}
